\chapter{Length, area, volume, mass}
\label{vol01:ch07}

\epigraph{The geometry of a body, and its mass, are the first quantities
an engineer can write down about it; everything else in the engineering
record builds on those four numbers.}{The working metrologist's rule,
phrased in the vocabulary of the SI Brochure \cite{std:bipm-si2019} and
the VIM \cite{std:jcgm-vim}.}

\chapmeta{Half-life: foundational. Archetypes: scaling (the geometric
quantities span twenty orders of magnitude); balance (mass conservation
in displacement and the residual-balance comparator).
Exercise target: 28. Project track: household + analysis (hybrid).
Hazard class: none.
Reader path default: standard, with sections explicitly tagged. Named
cases: Hubble Space Telescope primary mirror (1990), revisited from
chapters 2--4 as a wrong-ruler case in the failure section;
International prototype of the kilogram mass drift (1889--2019),
introduced in the mass section as an artefact-based-definition case
whose closure required the 2019 SI redefinition.}

\noindent
Length, area, volume, and mass are the first quantities an engineer
writes down about a physical object. The instruments that read them
range from a steel rule on a bench to a laser interferometer tied to
an optical frequency standard. Their arithmetic follows the
uncertainty rules of chapter 4, and their discipline is independent
checking: when two methods disagree outside their combined uncertainty,
the measurement process has found a wrong ruler.

\section{Length: ruler to interferometer}\pathtag{core}

Length is one of engineering's most common measurements, and its
practical range is unusually wide. A working engineer may meet lengths
from features in advanced integrated circuits, at a few nanometres, to
the orbital radius of a geostationary satellite, about $4.2 \times
10^{7} \,\si{\meter}$ from Earth's centre. No single instrument covers
that span. Length metrology is the catalogue of instruments matched to
scale, accuracy, material, and access.

\begin{archetype}[Scaling, geometric form]
Geometric quantities span enough decades that the reader needs an
internal ruler. A reader who carries a scale from $10^{-10} \,\si
{\meter}$ (atomic spacing) to $10^{7} \,\si{\meter}$ (continental
distance) has built seventeen decades of internal calibration. The
instruments below are organised by the decade they serve: each works
only over the few decades where its physical principle is matched to
the quantity's size.
\end{archetype}

\begin{figure}[p]
\centering
\begin{tikzpicture}[x=1cm, y=1cm, font=\small]

% Master horizontal axis: logarithmic decades of length, from 10^-10 to 10^7
\def\xmin{-0.4}
\def\xmax{17.4}

% Title
\node[anchor=south, font=\bfseries] at (8.5, 12.3) {Ladder of length instruments by decade of working range};

% Horizontal axis: decades of metres
\draw[->, very thick] (\xmin, 0) -- (\xmax, 0);
\foreach \i/\label in {0/$10^{-10}$, 1/$10^{-9}$, 2/$10^{-8}$, 3/$10^{-7}$, 4/$10^{-6}$, 5/$10^{-5}$, 6/$10^{-4}$, 7/$10^{-3}$, 8/$10^{-2}$, 9/$10^{-1}$, 10/$10^{0}$, 11/$10^{1}$, 12/$10^{2}$, 13/$10^{3}$, 14/$10^{4}$, 15/$10^{5}$, 16/$10^{6}$, 17/$10^{7}$} {
  \draw[thick] (\i, 0) -- (\i, -0.18);
  \node[anchor=north, font=\scriptsize] at (\i, -0.22) {\label};
}
\node[anchor=north west, font=\small] at (\xmax, -0.65) {metres};

% Reference decades, labelled
\node[anchor=south, font=\scriptsize\itshape] at (0.5, 0.05) {atom};
\node[anchor=south, font=\scriptsize\itshape] at (4.5, 0.05) {cell};
\node[anchor=south, font=\scriptsize\itshape] at (7.5, 0.05) {hair};
\node[anchor=south, font=\scriptsize\itshape] at (10.5, 0.05) {arm};
\node[anchor=south, font=\scriptsize\itshape] at (13.5, 0.05) {town};
\node[anchor=south, font=\scriptsize\itshape] at (16.5, 0.05) {country};

% Instruments as horizontal bars showing working range
% Each: \bar{name}{start}{end}{height}
\def\barh{0.5}
\newcommand{\instrbar}[5]{%
  \draw[ultra thick, #5] (#2, #1) -- (#3, #1);
  \node[anchor=west, font=\small] at (#3+0.15, #1) {#4};
}

\instrbar{1.0}{4}{6}{atomic-force microscope}{engblue}
\instrbar{1.7}{4}{8}{optical interferometer (laboratory)}{engblue}
\instrbar{2.4}{6}{8}{optical profilometer}{engred}
\instrbar{3.1}{6}{9}{coordinate measuring machine (precision)}{engred}
\instrbar{3.8}{7}{10}{micrometer}{engblack}
\instrbar{4.5}{7}{10}{Vernier caliper}{engblack}
\instrbar{5.2}{7}{11}{dial / digital indicator}{engblack}
\instrbar{5.9}{8}{12}{ruler and tape measure}{engblack}
\instrbar{6.6}{9}{13}{laser distance meter (hand-held)}{engred}
\instrbar{7.3}{11}{15}{total station (surveying)}{engred}
\instrbar{8.0}{12}{16}{GNSS receiver (positioning)}{engblue}
\instrbar{8.7}{14}{17}{satellite radar altimetry}{engblue}

% Reader-path band
\draw[densely dotted, thick] (7, -0.7) -- (10, -0.7);
\node[anchor=north, font=\scriptsize] at (8.5, -0.75) {everyday workshop band};

% Legend
\node[anchor=west, font=\scriptsize] at (0, 11.4) {\textbf{Legend:}};
\draw[ultra thick, engblack] (1.8, 11.4) -- (2.6, 11.4);
\node[anchor=west, font=\scriptsize] at (2.7, 11.4) {hand instrument};
\draw[ultra thick, engred] (5.5, 11.4) -- (6.3, 11.4);
\node[anchor=west, font=\scriptsize] at (6.4, 11.4) {laboratory or surveying};
\draw[ultra thick, engblue] (9.4, 11.4) -- (10.2, 11.4);
\node[anchor=west, font=\scriptsize] at (10.3, 11.4) {primary or wide-area};

\end{tikzpicture}
\caption[Length instruments organised by decade of working
range]{Length instruments organised by decade of working range, from
$10^{-10}$ to $10^{7} \,\si{\meter}$. Each bar marks the range over
which the instrument's physical principle is matched to the length
being measured; the overlaps between bars are what let a measurement
at one scale be cross-checked against a measurement at a neighbouring
scale. A working engineer's everyday band, marked beneath the axis,
covers about four decades from millimetre to metre. Source: editors'
synthesis of NIST and BIPM instrument-class descriptions.}
\label{fig:vol01:ch07:length-ladder}
\end{figure}


\Cref{fig:vol01:ch07:length-ladder} arranges the working length
instruments by the decade their physical principle covers. A reader
who develops a quick mental map between an instrument and its decade
has built one of the most useful internal indexes in engineering
metrology: the question "what instrument should I reach for?"
collapses to "which decade is my measurand in, and which bar in the
figure crosses that decade?"

\subsection{Working instruments by scale}

Seven instruments cover much of routine length work. Each occupies a
working range of two or three decades; the overlap among them is what
lets a measurement at one scale be checked against a measurement at a
neighbouring scale.

\begin{itemize}[leftmargin=*]
  \item \engterm{The ruler and tape measure}. Resolution at the
    millimetre, working range from a few centimetres to several
    metres. Verification against a traceable length standard at a
    stated reference temperature is the standard quality-control step;
    tape measures also need controlled tension and support. Tape
    measures stretch under tension and read short under compression;
    the working engineer learns to read at consistent tension, or to
    record the tension used.
  \item \engterm{The Vernier caliper}. Display resolution commonly
    ranges from $0.02$ to $0.05 \,\si{\milli\meter}$ on a Vernier
    instrument and $0.01 \,\si{\milli\meter}$ on a digital one, with
    working range often $150$ to $200 \,\si{\milli\meter}$ on a bench
    unit. Accuracy is worse than display resolution and must be taken
    from the instrument specification. The instrument's accuracy also
    depends on the operator's force on the jaws (a soft sample
    compresses), the parallelism of the jaws, and thermal expansion
    of the workpiece.
  \item \engterm{The micrometer}. Resolution $0.01 \,\si{\milli
    \meter}$ on a standard machine-shop instrument, $0.001 \,\si
    {\milli\meter}$ on a precision unit, working range $25 \,\si
    {\milli\meter}$ per instrument. Calibration against gauge blocks
    is documented in chapter 5 of this volume; the micrometer's
    precision sits well inside the realisation chain we built in
    chapter 3.
  \item \engterm{The dial gauge and digital indicator}. Resolution
    $1 \,\si{\micro\meter}$ at the high end, working range a few
    millimetres. Used for comparison measurements: the difference
    between the workpiece and a reference, rather than the workpiece's
    absolute dimension.
  \item \engterm{The coordinate measuring machine (CMM)}. A motorised
    probe traces points on a workpiece and reports their three-
    dimensional coordinates relative to the machine's reference
    frame. Working range runs from benchtop machines below a metre
    to large industrial machines of several metres. The uncertainty
    is machine-class and environment dependent, commonly specified as
    a length-dependent maximum permissible error such as $A + B L$,
    with the units of $A$, $B$, and $L$ stated by the relevant CMM
    specification. A CMM is the standard machine-shop bridge between
    a single linear measurement and a fully dimensioned
    three-dimensional part.
  \item \engterm{The laser distance meter and total station}. A
    pulsed or modulated laser is reflected from a target; the
    transit time or phase yields the distance. Hand-held units cover
    a few decimetres to roughly a hundred metres, with accuracy set
    by the instrument class and target. Surveying total stations can
    work over kilometre distances, but their specifications are
    usually given as millimetres plus parts per million under stated
    atmospheric and target conditions.
  \item \engterm{The interferometer}. A length is measured as a
    fraction of the wavelength of a stabilised light source; the
    wavelength is itself tied through the second's definition to the
    cesium-133 hyperfine transition \cite{std:bipm-si2019}.
    Frequency-stabilised laser interferometry is one route by which
    national metrology institutes realise length from the SI
    definition of the metre. In ordinary laboratory use, the
    uncertainty is usually dominated by alignment, air refractive
    index, temperature, and the mechanical setup, rather than by the
    optical frequency itself.
\end{itemize}

A working engineer rarely reaches beyond the calliper, micrometer,
and dial gauge in daily practice. The CMM, laser distance meter,
total station, and interferometer arrive when the application's
tolerance, range, or geometry demands them. The selection step is the
same one chapter 5 spelled out for sensors: state the measurand,
identify the candidate principles, compare specifications against the
tolerance, account for environment and integration, document the
choice.

\subsubsection*{Worked example: building a length reference from
gauge blocks}

The micrometer sits inside the realisation chain of chapter 3, but
the chain has a layer below it that the working engineer often
encounters directly: the \engterm{gauge block}. A gauge block is a
ceramic or steel slip ground to a length specified at $20 \,\si
{\degreeCelsius}$, with permissible deviations of length and
parallelism set by accuracy grade per ISO 3650 \cite{std:iso-3650}.
A reference length not available as a single block can be built up
by \engterm{wringing}: two flat, clean gauge-block faces pressed and
slid together adhere with a film of order nanometres between them,
and the stack's length is the sum of the individual lengths plus a
small residual film contribution.

Consider a stack of five blocks (50, 12, 5, 0.8, and $0.025 \,\si
{\milli\meter}$) that builds up a $67.825 \,\si{\milli\meter}$
reference. The standard uncertainty has two components. Each block
contributes its calibration uncertainty (taken from its certificate);
for a Grade-0 set, a representative standard uncertainty is
$0.05 \,\si{\micro\meter}$ per block in the size range $1$ to
$25 \,\si{\milli\meter}$ \cite{text:doiron-beers-2009}. With five
independently calibrated blocks, the combined calibration
contribution is $\sqrt{5} \times 0.05 \approx 0.11 \,\si{\micro
\meter}$. Each of the four wringing joints contributes a residual
film uncertainty; a conservative value is $0.02 \,\si{\micro\meter}$
per joint, so the wringing contribution is $\sqrt{4} \times 0.02 =
0.04 \,\si{\micro\meter}$. Combined in quadrature, the stack's
standard uncertainty is $\sqrt{0.11^{2} + 0.04^{2}} \approx 0.12 \,
\si{\micro\meter}$, dominated by the calibration term.
\repopath{volumes/01-quantity/ch07-length-area-volume-mass/code/gauge_block_stack.py}
encodes the calculation; the reader can substitute the certificates
of the actual blocks in use. The expanded uncertainty at $k = 2$ is
about $0.23 \,\si{\micro\meter}$ on a $67.825 \,\si{\milli\meter}$
stack, or roughly $3 \times 10^{-6}$ relative. The thermal-expansion
correction below changes whether that uncertainty is reportable.

\subsection{Length-specific failure modes}

Three error mechanisms recur across the length instruments above.

\engterm{Thermal expansion} biases every length measurement made on a
solid. A steel workpiece $1 \,\si{\meter}$ long, with linear thermal
expansion coefficient $\alpha \approx 1.2 \times 10^{-5} \,\si{\per
\kelvin}$, changes length by about $12 \,\si{\micro\meter}$ per
kelvin. A workshop $5 \,\si{\kelvin}$ warmer than the calibration
laboratory introduces a bias of $60 \,\si{\micro\meter}$ on a
metre-scale workpiece. The discipline is to record the workpiece
temperature, the instrument temperature, and the calibration
reference temperature. For dimensional specifications, correct to
the standard reference temperature of $20 \,\si{\degreeCelsius}$
where the specification requires it, or report the measurement at
its actual temperature with the correction stated
\cite{std:iso-1-2022}.

\engterm{Abbe error} arises when the measurement axis is offset from
the displacement axis of the instrument's reading scale. A small
angular error in the instrument's slide rotates the offset into a
length error proportional to the offset times the angle. Vernier
callipers are vulnerable; well-designed micrometers position the
spindle on the measurement axis and avoid the offset. The
Abbe principle, named for Ernst Abbe's 1890 statement of the
geometric requirement, is the design rule for high-accuracy
length instruments \cite{paper:bryan-abbe-1979}.

\begin{figure}[ht]
\centering
\begin{tikzpicture}[x=1cm, y=1cm, font=\small]

\node[anchor=south, font=\bfseries] at (7, 5.4) {Abbe error: offset between measurement and reading axes};

% Reading scale (top horizontal)
\draw[ultra thick, engblue] (1, 4.0) -- (13, 4.0);
\foreach \x in {1, 2, 3, 4, 5, 6, 7, 8, 9, 10, 11, 12, 13} {
  \draw[thick, engblue] (\x, 4.0) -- (\x, 4.15);
}
\node[anchor=south west, font=\scriptsize, engblue] at (13.1, 4.0) {reading scale axis};

% Measurement axis (lower horizontal, offset by h)
\draw[ultra thick, engred] (1, 2.0) -- (13, 2.0);
\node[anchor=south west, font=\scriptsize, engred] at (13.1, 2.0) {measurement axis};

% Vertical offset h
\draw[<->, thick] (0.5, 2.0) -- (0.5, 4.0) node[midway, left, font=\small] {$h$};

% Workpiece (rectangle along measurement axis)
\draw[ultra thick, fill=gray!20] (4, 1.7) rectangle (10, 2.3);
\node[anchor=south, font=\scriptsize] at (7, 2.32) {workpiece};

% Calliper jaws on workpiece (perfect alignment shown as overlay)
\draw[thick] (4, 4.0) -- (4, 2.0);
\draw[thick] (10, 4.0) -- (10, 2.0);

% Now show tilt: small angle theta exaggerated for clarity
% Tilted reading scale at angle theta about left jaw
\def\thetadeg{6}
\begin{scope}[shift={(4, 4.0)}, rotate=-\thetadeg]
  \draw[ultra thick, engblue, densely dashed] (0, 0) -- (9, 0);
  \draw[thick, engblue, densely dashed] (6, 0) -- (6, -0.15);
\end{scope}

% Annotation: angular error theta
\draw[thick, engblack] (4.6, 4.0) arc[start angle=0, end angle=-\thetadeg, radius=0.6];
\node[anchor=west, font=\small] at (4.7, 3.9) {$\theta$};

% Indicate length error: offset between true endpoint (10, 4) and tilted endpoint
\draw[->, thick, engred] (10, 4.2) -- (9.4, 3.4) node[anchor=west, font=\scriptsize] {reading shifts by $\approx h \theta$};

% Annotation: Abbe error formula
\node[anchor=north, font=\small] at (7, 1.2) {Abbe error: $\Delta L \approx h \cdot \theta$};
\node[anchor=north, font=\scriptsize] at (7, 0.7) {(when the measurement and reading axes are not coincident; here $h$ is the offset and $\theta$ is the small angular error)};
\node[anchor=north, font=\scriptsize] at (7, 0.15) {A well-designed micrometer puts the spindle on the measurement axis, $h \to 0$, eliminating the term.};

\end{tikzpicture}
\caption[Abbe error in length measurement]{An Abbe error arises when
the measurement axis (the line connecting the points whose distance
is being read) is offset from the reading-scale axis by a distance
$h$. A small angular tilt $\theta$ of the instrument's slide rotates
that offset into a length error of order $h \cdot \theta$. A Vernier
caliper has a finite offset and is vulnerable; a well-designed
micrometer puts the spindle on the measurement axis and drives the
offset to zero. The same principle underwrites the design of
high-accuracy coordinate measuring machines, where Abbe-compensating
scales and probe corrections are part of the published measurement
uncertainty model. Source: editors' schematic.}
\label{fig:vol01:ch07:abbe-error}
\end{figure}


\Cref{fig:vol01:ch07:abbe-error} illustrates the geometry. A
calliper with a $50 \,\si{\milli\meter}$ offset between its
measurement and reading axes, slipping by an angle of $0.01$ radian
($\approx 0.6^{\circ}$) during a measurement, produces an Abbe
contribution of about $500 \,\si{\micro\meter}$: half a millimetre,
many times the instrument's nominal resolution. The error has no
random signature in repeated readings under the same conditions,
because it tracks the geometric setup; only a measurement with a
different geometry, such as a micrometer on the same workpiece,
exposes it.

\engterm{Reference-frame ambiguity} matters when a CMM, laser
tracker, or surveying station establishes a coordinate frame that
the user must define. A CMM's reported coordinates are correct in
the machine's frame; whether they are correct in the workpiece's
design frame depends on how the user fixtured the workpiece and how
the alignment was reported. The discipline is to document the
fixture, the alignment, and the frame transformation alongside the
measurements.

\section{Area, volume, density}\pathtag{core}

Area and volume are derived geometric quantities. They follow from
length, but their measurement methods are not always reducible to a
chain of length measurements.

\subsection{Area}

For a regular planar shape, area follows from length:
$A = L W$ for a rectangle, $A = \pi R^{2}$ for a disc,
$A = \tfrac{1}{2} b h$ for a triangle. The uncertainty propagates by
the product-and-ratio rule of chapter 4: for $A = L W$ with
uncorrelated $L$ and $W$,

\[
\left(\frac{u(A)}{A}\right)^{2} = \left(\frac{u(L)}{L}\right)^{2} +
\left(\frac{u(W)}{W}\right)^{2}.
\]

For an irregular planar shape, the working methods are
\engterm{decomposition} (express the shape as a sum and difference of
regular subshapes whose areas add and subtract), \engterm{counting}
(overlay a grid of known cell area; count the cells inside the shape
plus a fraction for boundary cells), and \engterm{planimetry} (a
mechanical or optical instrument that integrates the boundary).
Counting is what most kitchen and workshop estimates do; planimetry
is the older, more precise method for tracing a planar boundary on a
drawing or photograph. A modern alternative is image analysis: a
photograph at a known scale is processed with software that counts
pixels inside the boundary.

\subsection{Volume}

Volume of a regular three-dimensional shape follows from lengths:
$V = L W H$ for a cuboid, $V = \pi R^{2} H$ for a cylinder, $V =
\tfrac{4}{3} \pi R^{3}$ for a sphere, $V = \tfrac{1}{3} A H$ for a
cone or pyramid. Uncertainty propagates by the same product-of-powers
rule, with the exponents on each length-input squared.

For an irregular shape, three working methods recur, and they form
the spine of this chapter's project.

\engterm{Water displacement (Archimedes)}. The object is submerged
in a measured volume of water; the volume of water displaced is the
object's volume. The reading is taken from the change in water level
in a graduated container, or from the mass of water displaced when
the object is suspended in a vessel that overflows. The method
assumes the object is fully wetted, that no air remains trapped on
the surface, and that the object is soluble in neither water nor
the surface coating used to wet it, and not absorbent. Its accuracy
is limited by the resolution of the graduation, the parallax error
in reading the meniscus, and the volume of water that adheres to the
object on removal \cite{hist:archimedes-floating}.

\begin{figure}[ht]
\centering
\begin{tikzpicture}[x=1cm, y=1cm, font=\small]

% Title
\node[anchor=south, font=\bfseries] at (5.5, 5.7) {Volume by water displacement (Archimedes)};

% Panel A: graduated cylinder before immersion
\begin{scope}[shift={(0,0)}]
  \node[anchor=south, font=\small\itshape] at (1.5, 5.0) {(a) before};
  % Cylinder body
  \draw[thick] (0.5, 0.4) -- (0.5, 4.8) (2.5, 0.4) -- (2.5, 4.8);
  \draw[thick] (0.5, 0.4) arc[start angle=180, end angle=360, x radius=1, y radius=0.25];
  \draw[thick, densely dashed] (0.5, 0.4) arc[start angle=180, end angle=0, x radius=1, y radius=0.25];
  % Water fill to level V0
  \fill[engblue!25] (0.5, 0.4) rectangle (2.5, 2.4);
  \draw[thick, engblue] (0.5, 2.4) -- (2.5, 2.4);
  \draw[thick, engblue] (0.5, 2.4) arc[start angle=180, end angle=0, x radius=1, y radius=0.18];
  % Graduations
  \foreach \y in {1.0, 1.5, 2.0, 2.5, 3.0, 3.5, 4.0} {
    \draw[thick] (2.4, \y) -- (2.5, \y);
  }
  % Level label
  \draw[->, thick, engred] (3.4, 2.4) -- (2.6, 2.4);
  \node[anchor=west, font=\small] at (3.4, 2.4) {$V_{0}$};
\end{scope}

% Panel B: graduated cylinder after immersion with stone
\begin{scope}[shift={(6,0)}]
  \node[anchor=south, font=\small\itshape] at (1.5, 5.0) {(b) after};
  % Cylinder
  \draw[thick] (0.5, 0.4) -- (0.5, 4.8) (2.5, 0.4) -- (2.5, 4.8);
  \draw[thick] (0.5, 0.4) arc[start angle=180, end angle=360, x radius=1, y radius=0.25];
  \draw[thick, densely dashed] (0.5, 0.4) arc[start angle=180, end angle=0, x radius=1, y radius=0.25];
  % Water to V1 (higher than V0)
  \fill[engblue!25] (0.5, 0.4) rectangle (2.5, 3.4);
  \draw[thick, engblue] (0.5, 3.4) -- (2.5, 3.4);
  \draw[thick, engblue] (0.5, 3.4) arc[start angle=180, end angle=0, x radius=1, y radius=0.18];
  % Stone immersed
  \fill[engblack!70] plot[smooth cycle, tension=0.7] coordinates {(1.0, 1.4) (1.4, 1.1) (2.0, 1.3) (2.1, 1.8) (1.8, 2.2) (1.3, 2.1) (0.9, 1.8)};
  % Graduations
  \foreach \y in {1.0, 1.5, 2.0, 2.5, 3.0, 3.5, 4.0} {
    \draw[thick] (2.4, \y) -- (2.5, \y);
  }
  % Level label
  \draw[->, thick, engred] (3.4, 3.4) -- (2.6, 3.4);
  \node[anchor=west, font=\small] at (3.4, 3.4) {$V_{1}$};
  % Stone label
  \node[anchor=west, font=\scriptsize] at (2.3, 1.6) {object};
\end{scope}

% Bottom equation
\node[anchor=north, font=\small] at (5.5, -0.1) {$V_{\text{object}} = V_{1} - V_{0}$};
\node[anchor=north, font=\scriptsize] at (5.5, -0.7) {valid only if the object is fully wetted, non-absorbent, and free of trapped air bubbles};

\end{tikzpicture}
\caption[Volume by water displacement]{Volume by water displacement.
The object's volume equals the rise in water level between the
before-immersion and after-immersion readings. The method assumes the
object is fully wetted, free of trapped air bubbles, and neither
soluble nor absorbent in water. The reading's uncertainty has three
sources: the cylinder's graduation interval (resolution), the
parallax of the meniscus reading, and the volume of water that
adheres to the object on removal. Section 7.5 returns to displacement
as one method in the chapter's three-method volume project. Source:
editors' schematic.}
\label{fig:vol01:ch07:volume-displacement}
\end{figure}


\Cref{fig:vol01:ch07:volume-displacement} shows the geometry. A
representative graduated cylinder with $1 \,\si{\milli\liter}$
graduations contributes a resolution-derived standard uncertainty of
$(1 / 2) / \sqrt{3} \approx 0.29 \,\si{\milli\liter}$ per reading
(uniform-distribution rule of chapter 4). A reading repeated five
times typically adds a Type A spread of a few tenths of a millilitre;
the two combine in quadrature.
\repopath{volumes/01-quantity/ch07-length-area-volume-mass/code/displacement_volume.py}
performs the combined-uncertainty calculation from a sequence of
readings.

\engterm{Geometric approximation}. The object is decomposed into a
sum (and difference) of simple shapes whose volumes the reader can
compute from length measurements. The accuracy depends on the fit
between the assumed shapes and the actual object: a smooth, nearly
spherical stone may be well-approximated by a sphere, while a stone
with sharp ridges and concave faces will produce a biased result if
either ridges or hollows are not accounted for.

\engterm{Image-based reconstruction}. Photographs from several
viewpoints, or a structured-light scan from a smartphone, produce a
three-dimensional model whose volume can be computed by software.
The accuracy depends on camera calibration, the number of viewpoints,
the quality of the surface texture, and the volume-extraction
algorithm.

A fourth method, \engterm{contact-based dimensional reconstruction},
uses a series of calliper or micrometer measurements at points
sampled across the object; software fits a parametric or mesh model
to those points and computes its volume. The method is slower than
image-based reconstruction but does not require visual texture, which
fails on uniform, transparent, or mirror-finish objects.

\paragraph{Worked example: 3D-scan-mesh volume integration.}
A handheld structured-light scanner produces a triangular surface
mesh of a cast aluminium bracket. The reader holds the volume
measurement to within $0.5\,\%$ for a density-based alloy
identification. The chain runs in four stages, each with its own
contribution to the volume uncertainty.

The scan stage delivers a point cloud at a stated lateral resolution
$\Delta_{\mathrm{lat}} \approx 0.1 \,\si{\milli\meter}$ and a stated
depth noise $\sigma_{z} \approx 0.05 \,\si{\milli\meter}$ at
$0.3 \,\si{\meter}$ working distance. For a bracket of nominal
dimensions $80 \times 60 \times 25 \,\si{\milli\meter}$ and nominal
volume $V_{0} \approx 60 \,\si{\milli\liter}$, the surface area is
of order $0.018 \,\si{\meter\squared}$. The depth-noise contribution
to the volume budget scales with $\sigma_{z} \times A_{\mathrm{surf}}$,
giving a per-scan uncertainty term of order $\sigma_{z} \cdot
A_{\mathrm{surf}} \approx 5 \times 10^{-5} \times 0.018 \approx
9 \times 10^{-7} \,\si{\meter\cubed} = 0.9 \,\si{\milli\liter}$,
or $1.5\,\%$ of $V_{0}$, before any averaging across redundant
viewpoints. Eight overlapping scans, registered together, average
the depth noise on the same surface patch by approximately
$\sqrt{8} \approx 2.8$, reducing the depth-noise contribution to
about $0.5\,\%$ of $V_{0}$.

The mesh-repair stage fills holes and snaps boundary edges to
produce a \engterm{watertight} mesh, the topological condition the
volume integrator needs. A non-watertight mesh has open edges that
the volume algorithm sees as ambiguous boundaries, and the result
can be off by tens of percent in either direction. The repair
discipline checks Euler characteristic ($V - E + F = 2$ for a closed
genus-zero surface, with $V$, $E$, $F$ the vertex, edge, and face
counts) and lists any remaining boundary loops before integration.
The repair stage introduces a discretisation contribution of order
the mesh edge length cubed times the number of newly inserted
triangles; for a bracket with $\sim 10\,000$ triangles and a
characteristic edge length of $1 \,\si{\milli\meter}$, the
repair-stage budget is of order $0.1\,\%$ of $V_{0}$.

The volume integration stage applies the divergence-theorem
identity. For a closed mesh with triangular facets, each triangle
$i$ with vertices $\vec{a}_{i}, \vec{b}_{i}, \vec{c}_{i}$ taken in
the outward-normal orientation, the enclosed volume is
\[
V = \frac{1}{6} \sum_{i} \vec{a}_{i} \cdot
(\vec{b}_{i} \times \vec{c}_{i}).
\]
The sum carries no truncation error for a closed mesh; the volume
the algorithm reports is the volume of the polyhedron the mesh
defines, exactly. The reported volume's error is the error in the
polyhedron's correspondence to the physical object, set by the
scan, mesh, and repair stages above. The
\repopath{volumes/01-quantity/ch07-length-area-volume-mass/code/mesh_volume.py}
script implements this integration on a Stanford-bunny test mesh
whose volume is independently catalogued, and reports $-0.08\,\%$
agreement with the catalogued value.

The thermal-expansion stage corrects the room-temperature scan to
the dimensional standard $20 \,\si{\degreeCelsius}$ of ISO 1
\cite{std:iso-1-2022}. Aluminium's linear-expansion coefficient is
$\alpha \approx 2.3 \times 10^{-5} \,\si{\per\kelvin}$; a scan at
$24 \,\si{\degreeCelsius}$ over-reports each linear dimension by
$\alpha \Delta T = 9.2 \times 10^{-5}$, and the volume by three
times that, $2.8 \times 10^{-4}$, or $0.028\,\%$ of $V_{0}$. The
correction is small but is recorded in the measurement file because
it depends on the assumed material (the same scan of a steel
bracket halves the correction; a polymer would multiply it).

Combining in quadrature: depth-noise $0.5\,\%$, mesh-repair
$0.1\,\%$, thermal $0.03\,\%$, plus a scale-calibration term
$0.2\,\%$ from the scanner's calibration certificate at the
$0.3 \,\si{\meter}$ working distance, gives a combined relative
standard uncertainty $u(V)/V \approx
\sqrt{0.5^{2} + 0.1^{2} + 0.03^{2} + 0.2^{2}} \approx 0.54\,\%$, a
hair above the $0.5\,\%$ project target. The mitigation reads off
the dominant term: add a ninth scan from the side that has only one
viewpoint, which by the $\sqrt{n}$ averaging reduces the depth-noise
contribution to $0.45\,\%$ and brings the budget under target. The
chain assumes a calibrated scanner against a traceable artefact (a
calibration sphere or a length bar), recorded in the scan file
metadata per ISO 10303-242
\cite{ch07-std:iso-iec-10303-242}; an uncalibrated scan reports a
volume but has no defensible uncertainty budget.

A counter-case is reported in the structured-light volume-measurement
literature. \cite{ch07-paper:zhang-3dscan-volume-2021} measure a
ceramic figurine of nominal volume $180 \,\si{\milli\liter}$ by both
water displacement and handheld scanning across mesh resolutions
$0.2$, $0.5$, and $1.0 \,\si{\milli\meter}$. The $0.2 \,\si{\milli
\meter}$ scan agrees with displacement to $0.3\,\%$; the $1.0 \,\si
{\milli\meter}$ scan reads $1.2\,\%$ low because the mesh
under-resolves a concave detail and the volume integrator treats the
missing concavity as solid rather than as void. The lesson for the
discipline: the mesh resolution must be set below the smallest
geometric feature whose volume contribution exceeds the project
tolerance, and the calibration check (compare against an independent
method on a known-shape artefact) is what verifies the resolution
choice for the part at hand.

\subsection{Density}

\engterm{Density} $\rho = m / V$ is the ratio of mass to volume. It
is the working bridge between this chapter's mass and volume sections
and the materials catalogue of Volume V. For an engineer trying to
identify a metal sample, density is often the cheapest first
discriminator: aluminium near $2700 \,\si{\kilogram\per\meter\cubed}$
distinguishes itself from steel near $7850 \,\si{\kilogram\per\meter
\cubed}$ from copper near $8960 \,\si{\kilogram\per\meter\cubed}$ from
lead near $11340 \,\si{\kilogram\per\meter\cubed}$ \cite{text:ashby2017}.

The uncertainty on density follows from the chapter 4 product-of-
powers rule applied to $\rho = m \cdot V^{-1}$:

\[
\left(\frac{u(\rho)}{\rho}\right)^{2} = \left(\frac{u(m)}{m}\right)^{2}
+ \left(\frac{u(V)}{V}\right)^{2},
\]

with the exponent on $V$ being $-1$, whose square is $1$. The larger
relative uncertainty usually dominates the density uncertainty; when
the two inputs are comparable, both contributions must be carried
in quadrature. The mass measurement is often the cheaper input to
refine: a $0.1 \,\si{\gram}$ scale gives a relative display step
near $10^{-3}$ on a hundred-gram sample, while a $1 \,\si{\gram}$
kitchen scale is nearer $10^{-2}$, comparable to the relative
volume uncertainty of a household graduated cylinder. The volume
measurement is therefore where the reader often spends precision
budget when density matters.

\begin{estimation}
\textbf{Estimate the density of a household object you can also
measure directly, then compare.}

\smallskip
\textit{Estimate, before reading on.}

\smallskip
A reader looks around and chooses a household object with a known
material. We pick a thick ceramic dinner plate. Published
engineering data put many dense ceramic bodies in the range of
roughly $2200$ to $2600 \,\si{\kilogram\per\meter\cubed}$, with the
exact value depending on clay body, porosity, and firing process
\cite{text:ashby2017}.

A working estimate proceeds without instruments. The plate is about
$25 \,\si{\centi\meter}$ in diameter, about $0.7 \,\si{\centi\meter}$
thick, and would feel about as heavy as half a litre of water. A disc
of those dimensions has volume $\pi (0.125)^{2} \times 0.007 \approx
3.4 \times 10^{-4} \,\si{\meter\cubed} = 340 \,\si{\milli\liter}$.
Estimating the mass at $750 \,\si{\gram}$ gives density
$750 / 340 \approx 2.2 \,\si{\gram\per\centi\meter\cubed}$, or
$2200 \,\si{\kilogram\per\meter\cubed}$. The estimate brackets the
published stoneware density; the factor-of-three test passes
comfortably.

As a worked comparison, take a measured plate with mass $830 \,\si
{\gram}$ and displaced volume $370 \,\si{\milli\liter}$. The
resulting density is $2240 \,\si{\kilogram\per\meter\cubed}$. The
estimate, $2200 \,\si{\kilogram\per\meter\cubed}$, is within about
$2\,\%$ of that measurement, much closer than the factor-of-three
target the estimation discipline asks for. The lesson is the
discipline, not the closeness: the estimate forced a sanity check
on the density-distinguishes-material claim above before any number
was looked up.
\end{estimation}

\subsection{Surface area: a separate discipline}

A reader sometimes needs the surface area of a body, distinct from
its planar projected area. Surface area enters heat-transfer,
reaction-rate, coating, and biological-transport calculations;
Volume IV and Volume V return to those uses with the governing
models. For a regular shape, surface area follows from
length by elementary geometry: $A_{\text{surface}} = 4\pi R^{2}$ for
a sphere, $A_{\text{surface}} = 2\pi R(R + H)$ for a closed cylinder.
For an irregular body, surface area is harder to measure than
volume; the standard methods are gas adsorption (BET, applied to
porous media) and image-based reconstruction with surface mesh
extraction. Volume V Chapter 7 takes up surface area in the materials
context; for this chapter's purposes, the reader should know that
surface area is a separate quantity from projected area and from
volume, and that conflating them is a recurring source of error in
heat-transfer and reaction calculations.

\section{Mass: balance, scale, recoil mass spectrometer}\pathtag{core}

The kilogram's definition changed in the 2019 SI redefinition, when
the Planck constant was assigned an exact numerical value
\cite{std:bipm-si2019,paper:stock-kibble-2018}. Chapter 3 covered the
realisation chain; this section covers the working instruments the
reader meets between primary mass realisation and the kitchen, and
returns at the end to the metrology-institute case (the
International prototype of the kilogram mass drift, 1889--2019) that
motivated the redefinition.

\subsection{The balance}

A \engterm{balance} compares an unknown mass against a reference
mass through a mechanical lever or a force-restoration servo. The
classical two-pan balance uses a beam pivoted at its centre; the
unknown sample on one pan is matched by reference masses on the
other until the beam reaches equilibrium. The method's accuracy is
set by the friction at the pivot, the linearity of the lever
geometry, the matching of the pans, and the reproducibility of
load placement.

\begin{archetype}[Balance, mass form]
The balance is the chapter's pure form of the balance archetype.
Equilibrium is the condition under which an unknown mass equals a
reference mass; the indication is whether the beam tilts. The
quantity reported is the reference mass that brings the indication
to zero. Later chapters apply the same pattern to charge balance,
momentum balance, energy balance, and information balance: the
indicator drives toward zero; the reference that achieves that
condition reports the unknown.
\end{archetype}

The modern \engterm{electronic balance} replaces the mechanical
beam with a magnetic-force restoration cell. A coil immersed in a
magnetic field generates a force opposing the weight on the pan; the
servo loop adjusts the coil current until the force restores
equilibrium; the current is reported as a mass after a calibration
that ties current-to-force to mass-to-weight at the local
gravitational acceleration. A laboratory analytical balance reaches
$0.1 \,\si{\milli\gram}$ resolution at $200 \,\si{\gram}$ full scale;
a microbalance reaches $0.1 \,\si{\micro\gram}$ at much smaller full
scales. Both are calibrated against mass references whose
traceability runs through the working laboratory chain we built in
chapter 3.

A classical equal-arm balance compares mass against mass. Local
gravitational acceleration cancels to first order because both pans
sit in the same gravitational field. A force-restoration electronic
balance instead measures the force needed to restore equilibrium
and reports mass after calibration. That instrument must be
calibrated for its site, because a move to a different latitude or
altitude changes the relation between mass and weight. Working
laboratory balances often include a built-in recalibration step
using an internal reference mass, run automatically on a
recalibration interval the operator sets.

\begin{figure}[ht]
\centering
\begin{tikzpicture}[x=1cm, y=1cm, font=\small]

% Title
\node[anchor=south, font=\bfseries] at (5, 5.6) {Equal-arm balance: mass against mass};

% Vertical column
\draw[ultra thick, engblack] (5, 0) -- (5, 4);

% Pivot at top
\fill[engblack] (5, 4) circle (0.08);
\node[anchor=west, font=\scriptsize] at (5.1, 4.05) {pivot};

% Horizontal beam, tilted slightly to indicate not-yet-balanced state
\draw[ultra thick, engblack] (1.6, 3.7) -- (8.4, 4.3);
% Reference mark for "level" indication
\draw[densely dotted] (1.4, 4) -- (8.6, 4);

% Left arm with pan
\draw[thick] (2, 3.78) -- (2, 3.2);
\draw[ultra thick, engblue, fill=engblue!20] (1.4, 3.0) rectangle (2.6, 3.2);
\node[anchor=north, font=\scriptsize] at (2, 2.95) {unknown sample};
\node[anchor=south, font=\small] at (2, 3.22) {$m$};

% Right arm with pan
\draw[thick] (8, 4.22) -- (8, 3.62);
\draw[ultra thick, engred, fill=engred!20] (7.4, 3.42) rectangle (8.6, 3.62);
\node[anchor=north, font=\scriptsize] at (8, 3.37) {reference masses};
\node[anchor=south, font=\small] at (8, 3.64) {$m_{\mathrm{ref}}$};

% Local gravity vector arrows (same on both sides, indicating cancellation)
\draw[->, thick, gray] (2, 2.3) -- (2, 1.5);
\node[anchor=west, font=\scriptsize, gray] at (2.1, 1.9) {$g$};
\draw[->, thick, gray] (8, 2.7) -- (8, 1.9);
\node[anchor=west, font=\scriptsize, gray] at (8.1, 2.3) {$g$};

% Equilibrium condition
\node[anchor=north, font=\small\itshape] at (5, 0.8) {Equilibrium: $m \cdot g = m_{\mathrm{ref}} \cdot g \implies m = m_{\mathrm{ref}}$};
\node[anchor=north, font=\scriptsize] at (5, 0.25) {$g$ cancels because both pans sit in the same gravitational field};

% Pointer below pivot
\draw[thick, engred] (5, 3.95) -- (5, 2.4);
\draw[thick] (4.5, 2.4) -- (5.5, 2.4);
\node[anchor=west, font=\scriptsize] at (5.55, 2.4) {indicator};

\end{tikzpicture}
\caption[Equal-arm balance schematic]{The equal-arm balance compares
an unknown mass on one pan against reference masses on the other. At
equilibrium, the beam's pointer reads zero; the gravitational
acceleration $g$ appears on both sides and cancels to first order, so
the reported quantity is mass, not weight. A force-restoration
electronic balance replaces the mechanical beam with a coil in a
magnetic field; the coil current at equilibrium is reported as a mass
after a site-specific calibration that ties current to weight at the
local value of $g$. Source: editors' schematic.}
\label{fig:vol01:ch07:equal-arm-balance}
\end{figure}


\Cref{fig:vol01:ch07:equal-arm-balance} shows the equal-arm pattern.
The geometric cancellation of $g$ is the structural reason a
two-pan balance reports mass in any laboratory on Earth's surface
without site-specific calibration. The force-restoration cell loses
that cancellation in exchange for a single pan, faster reading, and
electronic output; the calibration discipline it requires is the
price of those benefits.

\subsection{The scale}

A \engterm{scale} reports weight, often labelled in mass units. A
bathroom scale, a postal scale, a kitchen scale, and a truck scale
usually read the gravitational force on the load and report a
mass-equivalent value. The reading is tied to the gravitational
acceleration assumed during calibration. If a force scale
calibrated in northern Europe is used near the equator at high
altitude, the local value of $g$ can differ by several tenths of a
percent; without recalibration, the indicated mass shifts by the
same fraction.

Scales typically use a strain-gauge load cell: a metal flexure
deforms under load, strain gauges bonded to the flexure measure the
deformation, and the conditioning electronics convert the strain
reading to a calibrated mass. Chapter 5 of this volume covered strain
gauges as sensors; the load cell is the most common engineering
deployment of the principle.

\subsection{The mass spectrometer}

A \engterm{mass spectrometer} measures the mass-to-charge ratio of
ionised atoms or molecules. The dominant working principles include
quadrupole filtering (where ions of a chosen mass-to-charge are
selected by oscillating electric fields), magnetic-sector deflection
(where ions of different mass-to-charge follow different curved paths
in a magnetic field), and time-of-flight separation (where ions
accelerated through a fixed voltage drift across a known distance,
with mass inferred from the transit time).

Mass spectrometry belongs to chemistry, biochemistry, materials
analysis, and nuclear science. It reports mass-to-charge ratio for
atomic and molecular ions, and its calibration is a specialist
procedure. It belongs in the catalogue because mass work ranges from
ship-scale weighbridges, on the order of $10^{8}$ to $10^{9} \,\si
{\kilogram}$ for the largest loaded vessels, to atomic-scale
measurements in daltons. Volume V takes up mass spectrometry in the
materials-analysis context.

\subsection{The Kibble balance and the post-2019 kilogram}

A working laboratory balance reports mass against a reference whose
own provenance traces, by a chain of calibrations, to a primary
realisation of the kilogram. Until 2019 the primary realisation was
the international prototype of the kilogram (IPK), a single
platinum-iridium artefact at the BIPM. The 2019 SI redefinition
replaced that artefact with a fixed numerical value of the Planck
constant, and the modern primary realisation is the
\engterm{Kibble balance}, an electrical-mechanical instrument that
ties mass to the electrical units (volt and ohm) through the
Josephson and von Klitzing constants \cite{web:nist-kibble-balance,%
paper:stock-kibble-2018,web:bipm-mise-en-pratique-kilogram-2019}.

\begin{figure}[ht]
\centering
\begin{tikzpicture}[x=1cm, y=1cm, font=\small]

% Title
\node[anchor=south, font=\bfseries] at (6, 6.3) {Kibble balance: kilogram realisation from the Planck constant};

% ====== Weighing mode (left) ======
\begin{scope}[shift={(0, 0)}]
  \node[anchor=south, font=\small\bfseries] at (2.5, 5.5) {(a) Weighing mode};

  % Mass pan
  \draw[thick, fill=engred!20] (1.5, 3.8) rectangle (3.5, 4.0);
  \node[anchor=south, font=\small] at (2.5, 4.05) {test mass $m$};

  % String / suspension
  \draw[thick] (2.5, 3.8) -- (2.5, 2.6);

  % Coil in magnetic field
  \draw[thick, engblue, fill=engblue!15] (1.8, 2.2) rectangle (3.2, 2.6);
  \node[anchor=west, font=\scriptsize] at (3.3, 2.4) {coil, current $I$};

  % Magnet poles
  \draw[thick, fill=gray!30] (0.6, 2.0) rectangle (1.6, 2.8);
  \draw[thick, fill=gray!30] (3.4, 2.0) rectangle (4.4, 2.8);
  \node[anchor=south, font=\scriptsize] at (1.1, 2.8) {N};
  \node[anchor=south, font=\scriptsize] at (3.9, 2.8) {S};

  % Field lines (radial)
  \foreach \y in {2.15, 2.4, 2.65} {
    \draw[->, gray, thick] (1.7, \y) -- (3.3, \y);
  }
  \node[anchor=north, font=\scriptsize, gray] at (2.5, 2.0) {field $B$};

  % Gravity arrow
  \draw[->, thick, engblack] (5.0, 4.0) -- (5.0, 3.2);
  \node[anchor=west, font=\small] at (5.05, 3.6) {$mg$};

  % Equation
  \node[anchor=north, font=\small] at (2.5, 1.3) {$mg = B L I$};
  \node[anchor=north, font=\scriptsize] at (2.5, 0.75) {coil current balances weight};
\end{scope}

% ====== Velocity mode (right) ======
\begin{scope}[shift={(7.5, 0)}]
  \node[anchor=south, font=\small\bfseries] at (2.5, 5.5) {(b) Velocity mode};

  % Empty pan
  \draw[thick] (1.5, 3.8) rectangle (3.5, 4.0);
  \node[anchor=south, font=\scriptsize] at (2.5, 4.05) {pan empty};

  % Suspension
  \draw[thick] (2.5, 3.8) -- (2.5, 2.6);

  % Coil
  \draw[thick, engblue, fill=engblue!15] (1.8, 2.2) rectangle (3.2, 2.6);
  \node[anchor=west, font=\scriptsize] at (3.3, 2.4) {coil, voltage $U$};

  % Magnet poles
  \draw[thick, fill=gray!30] (0.6, 2.0) rectangle (1.6, 2.8);
  \draw[thick, fill=gray!30] (3.4, 2.0) rectangle (4.4, 2.8);
  \node[anchor=south, font=\scriptsize] at (1.1, 2.8) {N};
  \node[anchor=south, font=\scriptsize] at (3.9, 2.8) {S};
  \foreach \y in {2.15, 2.4, 2.65} {
    \draw[->, gray, thick] (1.7, \y) -- (3.3, \y);
  }

  % Velocity arrow
  \draw[->, thick, engred] (2.5, 2.7) -- (2.5, 3.5);
  \node[anchor=west, font=\small] at (2.55, 3.1) {$v$};

  % Equation
  \node[anchor=north, font=\small] at (2.5, 1.3) {$U = B L v$};
  \node[anchor=north, font=\scriptsize] at (2.5, 0.75) {moving coil induces voltage};
\end{scope}

% Combination equation at the bottom
\draw[->, thick] (4.6, 0.0) -- (6.6, 0.0);
\node[anchor=north, font=\small] at (6, -0.4) {combine};
\node[anchor=north, font=\small] at (6, -1.0) {$mg \cdot v = U \cdot I \implies m = \dfrac{U I}{g v}$};
\node[anchor=north, font=\scriptsize] at (6, -1.7) {$BL$ cancels; with $h$ fixed in the 2019 SI, $U/I$ ties electrical to mechanical units};

\end{tikzpicture}
\caption[Kibble balance schematic]{The Kibble balance realises the
kilogram from the Planck constant in two phases. In weighing mode
(panel a), the current $I$ in a coil placed in a radial magnetic
field $B$ is adjusted so the magnetic force on the coil equals the
weight of the test mass: $mg = BLI$. In velocity mode (panel b), the
same coil moves through the field at velocity $v$, inducing a voltage
$U = BLv$. Multiplying the two equations cancels the geometric factor
$BL$ and gives $m = UI/(gv)$. The 2019 SI redefinition fixed the
Planck constant $h$, which ties electrical units (volt and ohm via
the Josephson and von~Klitzing constants) to mechanical units, so the
right-hand side becomes an SI-traceable expression for $m$. Source:
editors' schematic following the NIST Kibble-balance description
\cite{web:nist-kibble-balance,std:bipm-si2019}.}
\label{fig:vol01:ch07:kibble-balance}
\end{figure}


\Cref{fig:vol01:ch07:kibble-balance} shows the two-mode operation.
In weighing mode, the magnetic force on a current-carrying coil in
a radial field opposes the weight of the test mass; in velocity
mode, the same coil moved at known velocity through the same field
induces a voltage. Multiplying the two equations cancels the
geometric factor $BL$ and yields $m = U I / (g v)$, with $U$, $I$,
$g$, and $v$ all SI-traceable. The 2019 redefinition closed the
loop: with the Planck constant fixed, the electrical-to-mechanical
conversion is exact, and the Kibble balance realises the kilogram at
a relative uncertainty currently of order $10^{-8}$
\cite{paper:stock-kibble-2018}. The X-ray crystal density (XRCD)
method, which counts silicon atoms in a polished sphere of known
crystal structure, provides an independent primary realisation; the
agreement between the two methods is the cross-check that lets the
SI tolerate either realisation as primary.

\subsection{Failure case: the International prototype of the
kilogram, 1889--2019}

The reason an artefact-based mass definition was retired in 2019 is
a 130-year case of unfalsifiable drift. The International prototype
of the kilogram mass drift (1889--2019), a platinum-iridium artefact
sanctioned by the 1st CGPM in 1889 and held under nested bell jars
at the BIPM at S{\`e}vres, was the world's reference mass for over
a century. Periodic verifications against the BIPM's official
witness copies and the national prototypes distributed to member
states showed apparent drifts of order tens of micrograms among the
artefacts over decades; the third verification of 1988--1992
quantified the divergence at the standard-deviation level of several
tens of micrograms across the cohort \cite{paper:girard-ipk-1994,%
web:bipm-kilogram-history}.

\paragraph{Technical mechanism.}
Candidate physical mechanisms for the drift include adsorption and
desorption of contaminants on the platinum-iridium surfaces (gain
and loss of several micrograms per cleaning cycle had been
demonstrated), diffusion of hydrogen and other gases into the bulk
metal, and mechanical wear from periodic handling. The surface was
finished to optical quality, and the artefact's $39 \,\si
{\milli\meter}$ dimensions had been chosen in 1889 to minimise
surface area for a given volume; even so, the contamination layer
that the standardised washing protocol removed was of order
nanometres thick, and its mass scaled with surface area in
micrograms. The crucial measurement asymmetry was definitional. The
IPK was, by definition, exactly one kilogram, so its absolute drift
could not be measured: any divergence from a witness copy was
reportable as the copy's drift, not the IPK's. The cohort scatter
was real and quantifiable; the assignment of that scatter to one
artefact or another was not \cite{paper:girard-ipk-1994}.

\paragraph{Organisational context.}
The IPK governance structure made the artefact's primacy
self-fulfilling. By the 1875 Metre Convention and the 1889 CGPM
sanction, the IPK was, by definition, the unit of mass; no
procedural mechanism existed to declare a definitional change in the
SI kilogram on the basis of a verification result. The path out of
the impasse required two decades of metrological research into
electrical analogues of mass, originating with Bryan Kibble's NPL
proposal in 1975, plus independent realisations by XRCD with isotopically
purified silicon spheres. Once multiple methods agreed at the
$10^{-8}$ relative-uncertainty level, the CGPM resolutions of 2011
and 2018 approved the redefinition, and the 2019 SI Brochure
recorded the change \cite{paper:stock-kibble-2018,std:bipm-si2019,%
web:bipm-mise-en-pratique-kilogram-2019}.

\paragraph{Lesson for mass realisation.}
An artefact-based unit definition has the same structural blind spot
as the Hubble RNC discussed below: the reference instrument's own
drift cannot be measured against itself. The IPK case is the
slower-motion, larger-stakes version. The closing of the loop took
130 years because the alternative required a different physical
principle to be developed to the required uncertainty. The
operational lesson for a working engineer is structural: any
measurement traced to a single physical artefact, however well
maintained, has an unmeasured drift component whose magnitude is
bounded only by the agreement among nominally identical copies of
that artefact. The discipline that retired the IPK is the same
cross-checking discipline of section 7.5, applied at the highest
metrological level: when the cohort of supposedly identical
references diverges, the unit must be redefined against something
that does not depend on the artefact.

\subsection{Mass-specific failure modes}

Three failure modes recur in mass measurements.

\engterm{Air buoyancy}. Air has density. A sample in air loses
apparent weight by the weight of the air it displaces, and the
apparent mass correction becomes significant for low-density samples
or high-precision weighing; it also depends on the density of the
reference weights used in calibration. A polystyrene foam sample of volume
$10^{-3} \,\si{\meter\cubed}$ displaces about $1.2 \,\si{\gram}$ of
air at sea level; a balance calibrated against denser reference
masses reads about $1.2 \,\si{\gram}$ less than the polystyrene's
true mass. For most kitchen and workshop work, the correction is
negligible. For analytical balance work on low-density samples or
for high-precision work on any sample, the correction is part of the
measurement \cite{text:taylor1997}.

\engterm{Electrostatic charge}. A dry sample on an analytical balance
in a low-humidity room may carry a static charge sufficient to
attract or repel the pan. The reading drifts as the charge equilibrates
or as the sample is displaced by image charges in the pan. The
mitigation is to ground the pan, raise the humidity, or use a
charge-dissipation device near the balance.

\engterm{Mechanical settling}. A balance and a scale both have a
settling time after a load change. A reading taken before the
instrument has settled will drift; a reading taken after the
instrument has settled but on a vibrating surface will fluctuate. The
discipline is to wait for the indicator to stabilise and to record
either the stable indication or the average and spread of the
fluctuations, with conditions noted.

\begin{estimation}
\textbf{Estimate the relative uncertainty of a bathroom scale
reading on a representative adult, then test the limits of what the
instrument can defensibly resolve.}

\smallskip
\textit{Estimate, before reading on.}

\smallskip
A reader picks a household bathroom scale and an adult of mass near
$70 \,\si{\kilogram}$. The scale typically reports in $0.1 \,\si
{\kilogram}$ steps; the resolution-derived standard uncertainty per
reading is $(0.1 / 2)/\sqrt{3} \approx 0.029 \,\si{\kilogram}$, or
about $0.04\,\%$ of the load. The repeatability spread of a
consumer scale, from repeated stepping on and off in a single
session, is typically a few tenths of a kilogram; let us estimate
$\pm 0.2 \,\si{\kilogram}$ standard deviation for a representative
unit, or $0.3\,\%$ of the load. Calibration drift between purchase
and a year of use, also a few tenths of a kilogram on a typical
strain-gauge load cell, contributes a similar Type B component.
Combining the three in quadrature gives a total relative standard
uncertainty of order $0.4\,\%$, with the repeatability and the
drift dominating the resolution.

\smallskip
The estimation has an operational consequence. A reader who steps
on the scale at the beginning and the end of a week and reads a
$0.5 \,\si{\kilogram}$ change is reading a quantity within the
scale's combined standard uncertainty: the change is not
statistically distinguishable from zero on a single-reading basis.
A reader who averages five readings per day for the week reduces
the Type A component by $\sqrt{5}$ but does not change the
calibration drift, leaving the latter as the floor. To detect a
genuine $0.5 \,\si{\kilogram}$ change at the $r \leq 2$ threshold
of section 7.5, the protocol would need either a higher-class
scale, daily averaging across many readings, or a different
quantity (body-water content, food intake records) measured
independently.

\smallskip
The same arithmetic explains why a laboratory analytical balance
with $0.1 \,\si{\milli\gram}$ resolution at $200 \,\si{\gram}$ full
scale is not a different instrument from the bathroom scale by a
small factor: it is a different instrument by four orders of
magnitude, and its applications, costs, and calibration disciplines
scale accordingly.
\end{estimation}

\section{Specialised quantities: viscosity, hardness, surface
roughness}\pathtag{standard}

Three related quantities recur often enough to name here before
Volume V treats them in detail. Each is at the recognition level:
the working sensor, the working instrument, and the failure mode
the reader should know about.

\subsection{Viscosity}

\engterm{Viscosity} is a fluid's resistance to shear deformation. The
SI unit of dynamic viscosity is the pascal second, $\si{\pascal\second}$;
the older unit, the poise, equals $0.1 \,\si{\pascal\second}$. Water
at $20 \,\si{\degreeCelsius}$ has dynamic viscosity near
$1 \,\si{\milli\pascal\second}$. Lubricating oils and honey occupy
much higher, strongly temperature-dependent ranges, so any quoted
value must carry temperature, grade, and source.

The working instruments measure viscosity by the response of a probe
to fluid drag. A \engterm{rotational viscometer} drives a spindle in
the fluid at known angular velocity and measures the torque required
to maintain the velocity; the ratio of torque to velocity, divided
by a geometric factor, yields the viscosity. A \engterm{falling-ball
viscometer} drops a calibrated sphere through the fluid in a vertical
tube; the terminal velocity, after the ball has reached steady state,
gives the viscosity through Stokes's law for laminar flow. A
\engterm{capillary viscometer} drains a known volume of fluid through
a narrow tube under gravity; the time required gives the kinematic
viscosity through the Hagen-Poiseuille relation. Each instrument has
a working range, a temperature requirement (viscosity is strongly
temperature-dependent), and a calibration reference (a reference
fluid of known viscosity).

The recurring failure mode is temperature. Viscosity is strongly
temperature-dependent; for many liquids, a modest temperature
change moves the reading by engineering-significant amounts. A
viscosity measurement at an unrecorded temperature is incomplete.
The discipline is to thermostat the sample, record the temperature,
and report viscosity at a stated reference temperature.

\subsection{Hardness}

\engterm{Hardness} is a material's resistance to localised plastic
deformation. The working measurement methods are indentation tests:
an indenter of a defined geometry is pressed into the material under
a defined load, and the size of the residual indentation is
measured. The dominant working scales are Brinell (a hardened steel
or tungsten carbide ball, indenting a flat surface, with the
indentation diameter measured), Vickers (a square-pyramid diamond
indenter, with the diagonal of the residual square measured), and
Rockwell (a conical or spherical indenter, with the depth of
indentation reported on a relative scale). Each scale has a defined
load, indenter geometry, and conversion convention; the scales are
not directly interchangeable, and a hardness reported on one scale
must be converted to another only through tables or correlations
specific to the material class.

The recurring failure mode is conversion. A reader who reads a
"hardness of 50" without the scale qualifier, such as HRC, HV, or
HB, cannot interpret the number. Converting among hardness scales
is material-dependent and approximate; use a standard or a
material-specific correlation table, and report the source with the
converted value. A hardness number reported without its scale is
incomplete in the same way a length reported without its unit is
incomplete \cite{text:ashby2017}.

\paragraph{Worked example: Vickers depth-of-case profile on a
carburised gear tooth.}
A carburised steel gear tooth is specified with a case-hardening
depth (CHD) of $0.80 \pm 0.10 \,\si{\milli\meter}$. The depth is
defined per ISO 2639 as the perpendicular distance from the
surface to the point where the Vickers hardness HV1 (test load
$1 \,\text{kgf}$) crosses $550$, the conventional
case-to-core threshold \cite{ch07-std:iso-2639-2002}. The reader
runs a Vickers microhardness profile to verify the specification.

The instrument is a microhardness tester with a Vickers diamond
indenter, calibrated against a reference block at HV1 $= 700$ and
HV1 $= 200$. The Vickers number is computed from the mean
indentation diagonal $d$ by
\[
\mathrm{HV} = \frac{2 F \sin(136^{\circ}/2)}{d^{2}} \approx
\frac{1.8544 \, F}{d^{2}},
\]
with $F$ in newtons and $d$ in millimetres returning HV in
units of $\text{kgf}/\text{mm}^{2}$ by
convention \cite{ch07-std:iso-6507-1-2023}. At HV1 $= 550$, the
expected diagonal is $d = \sqrt{1.8544 \times 9.807 / 550} \approx
0.182 \,\si{\milli\meter}$, or $182 \,\si{\micro\meter}$. The
operator measures $d$ to the nearest $0.5 \,\si{\micro\meter}$
under a calibrated optical reticle; the reading uncertainty is
$\sigma_{d} \approx 0.5 \,\si{\micro\meter}$ uniform, contributing
a hardness uncertainty $\sigma_{\mathrm{HV}} / \mathrm{HV} =
2 \sigma_{d} / d \approx 0.5\,\%$ per measurement.

The indentation spacing requirement is the working depth-of-case
constraint. ISO 6507-1 requires adjacent indentations be separated
by at least $2.5 d$ centre-to-centre to avoid work-hardened-zone
overlap \cite{ch07-std:iso-6507-1-2023}. At $d = 0.18 \,\si{\milli
\meter}$ the minimum spacing is $0.45 \,\si{\milli\meter}$,
coarser than the $0.10 \,\si{\milli\meter}$ tolerance on CHD that
the specification demands. A staggered double-row layout (two
parallel rows offset by half the spacing) doubles the achievable
depth resolution to $0.22 \,\si{\milli\meter}$, still too coarse.
The defensible procedure runs HV0.1 (test load $0.1 \,\text{kgf}$)
instead, with expected diagonal $d \approx 57 \,\si
{\micro\meter}$ at HV $= 550$, minimum spacing $0.14 \,\si{\milli
\meter}$, and resolution adequate for the specification
\cite{ch07-std:astm-e384-2022}. The trade-off is a higher
reading-uncertainty term ($\sigma_{d} / d \approx 1\,\%$ instead of
$0.3\,\%$ at the smaller indent), and a reduced depth-of-influence
in the material below the indent, which for HV0.1 is of order
$5 d \approx 0.3 \,\si{\milli\meter}$ rather than $0.9 \,\si{\milli
\meter}$.

The reader runs 16 HV0.1 indents at $0.10 \,\si{\milli\meter}$
depth intervals from $0.10 \,\si{\milli\meter}$ to $1.60 \,\si
{\milli\meter}$ below the surface. The reported profile (idealised
exponential decay typical of carburised steel) is
$\mathrm{HV}(z) \approx \mathrm{HV}_{\mathrm{core}} + (\mathrm{HV}
_{\mathrm{surf}} - \mathrm{HV}_{\mathrm{core}})
\exp(-z / \lambda)$, with $\mathrm{HV}_{\mathrm{core}} = 250$,
$\mathrm{HV}_{\mathrm{surf}} = 720$, and characteristic length
$\lambda \approx 0.45 \,\si{\milli\meter}$ for this carburising
recipe. Solving $\mathrm{HV}(z_{*}) = 550$ for $z_{*}$ gives
$z_{*} = -\lambda \ln((550 - 250)/(720 - 250)) = -0.45 \times
\ln(300/470) \approx 0.20 \,\si{\milli\meter}$. The measured CHD,
read off the profile by the same logic at the points immediately
above and below HV1 $= 550$, is the depth at which the smoothed
HV0.1 curve crosses the threshold.

The uncertainty budget on $z_{*}$ runs through three terms. The
HV-reading term at the threshold contributes
$\sigma_{z, \mathrm{HV}} = \lambda \, \sigma_{\mathrm{HV}} /
(\mathrm{HV}_{\mathrm{surf}} - \mathrm{HV}_{\mathrm{core}}) \cdot
\exp(-z_{*}/\lambda) \approx 0.45 \times 5 / 470 \times 0.64
\approx 0.003 \,\si{\milli\meter}$ at the threshold (negligible).
The depth-positioning term, set by the metallographic-mount
sectioning and the stage-positioning uncertainty, is $\sigma_{z,
\mathrm{pos}} \approx 0.02 \,\si{\milli\meter}$. The HV1-to-HV0.1
load-conversion term, which the standard
\cite{ch07-std:astm-e384-2022} flags as nontrivial for the case
threshold, is bounded experimentally by running HV1 indents at the
coarser spacing on the same sample after the HV0.1 profile is
complete; the offset between HV0.1 and HV1 at the threshold is of
order $20 \,\text{HV}$ on this steel class
\cite{ch07-paper:vickers-case-depth-2018}, contributing
$\sigma_{z, \mathrm{conv}} \approx \lambda \times 20 / 470 \approx
0.02 \,\si{\milli\meter}$. Combined in quadrature, $u(z_{*})
\approx 0.03 \,\si{\milli\meter}$, well inside the
$0.10 \,\si{\milli\meter}$ specification tolerance.

The discipline restated. A bare HV reading is incomplete without
its load specification (HV1, HV10, HV0.1) and without the spacing
and depth-of-influence justifications that the chosen load must
satisfy. The reader who reports "CHD $= 0.80 \,\si{\milli\meter}$
by Vickers profile" without naming the load class is reporting a
number whose uncertainty is unknowable, because the conversion
between Vickers loads at the case threshold is material-specific
and adds a contribution that the bare number cannot expose.

\subsection{Surface roughness}

\engterm{Surface roughness} is the deviation of a surface from its
ideal smooth form, measured at a length scale below the surface's
intended geometric features. The working units are the arithmetic
mean roughness $R_{a}$ and the root-mean-square roughness $R_{q}$,
both reported in micrometres or nanometres depending on the surface.
Machined and polished surfaces span orders of magnitude in $R_{a}$,
from micrometre-scale roughness for many machined finishes to
nanometre-scale roughness for high-quality optical polishing. Quote
a process-specific value only with the measurement method and
source.

The working instruments are the contact stylus profilometer (a
stylus traces the surface, recording its vertical displacement
versus position; the resulting profile yields $R_{a}$ and $R_{q}$),
the optical profilometer (interferometric or confocal techniques
sample the surface without contact, suitable for delicate or soft
samples), and the atomic force microscope (a sharp tip scans the
surface at the nanometre scale, suitable for the smoothest finishes).

The recurring failure mode is sampling length and filtering. A
profilometer reading is sensitive to the length over which the
profile is sampled and to the spatial-frequency filter used to
separate roughness from waviness and form. Two laboratories
measuring "the same surface" can report different $R_{a}$ values if
their evaluation lengths and filters differ. The discipline is to
specify the evaluation length, cutoff or nesting index, filter
type, and governing standard alongside the roughness value
\cite{std:iso-21920-2-2021}.

\section{Cross-checking: independent measurements of the same
quantity}\pathtag{core}

The practical test is independent measurement of the same quantity.
A single measurement reports a value with an uncertainty derived
from chapter 4's algebraic rules. A second, independent measurement
reports another value with another uncertainty. If the two agree
within their combined uncertainty, confidence in the quantity grows.
If they disagree by more than the combined uncertainty, at least one
method has an uncharacterised error source.

Cross-checking is the most direct wrong-ruler detector available
inside a local measurement task. An internal uncertainty budget
reports what the measurement chain believes about itself; an
independent measurement reports what nature delivers when asked
through a different chain. Disagreement between the two is the
diagnostic; agreement is the credential.

\subsection{The discipline}

The procedure has four steps.

First, choose two or more methods whose physical principles, error
sources, and instruments are genuinely independent. A length measured
with two callipers from the same drawer of the same workshop is not
two independent measurements; both inherit the workshop's
temperature, the operator's parallax, and the calibration history of
the same calibration laboratory. A length measured with a calliper
and with a laser distance meter is closer to independent: the
parallax error of the calliper has no analogue in the laser
instrument, and the laser's distance reference is not in the
calliper's chain.

Second, compute each method's value and combined standard
uncertainty using the chapter 4 rules. The two reports become a pair
$(x_{A}, u_{A})$ and $(x_{B}, u_{B})$.

Third, compute the difference and its uncertainty. For uncorrelated
methods, the difference $\Delta = x_{A} - x_{B}$ has standard
uncertainty $u(\Delta) = \sqrt{u_{A}^{2} + u_{B}^{2}}$. The
\engterm{compatibility ratio} is

\[
r = \frac{|\Delta|}{u(\Delta)}.
\]

A value of $r$ near $1$ or below means the two measurements agree
within their combined standard uncertainty. A value near $2$ is a
warning threshold: the methods disagree by about twice their
combined standard uncertainty, so the application must decide
whether to investigate, widen the uncertainty budget, or repeat the
measurement.

Fourth, if $r$ exceeds the threshold the application requires, look
for a previously uncharacterised error source. Either an
uncertainty was understated or a systematic effect was missed. The
investigation tightens the budget; the next round of measurements
either confirms the agreement or reveals a deeper effect.

\subsection{What independence buys}

Independent measurements catch error sources that the methods do not
share. A length measured with both a calliper and a laser distance
meter detects most measurement-process errors specific to either
instrument, since the calliper's parallax has no analogue in the
laser, and the laser's modulation reference is not in the calliper's
chain. A bias common to both methods, such as a workshop temperature
shift acting on the steel workpiece, propagates through both readings
together and goes undetected. Independence is the catalogue of error
sources the two methods do, and do not, share.

Three error families propagate through cross-checking differently.
\engterm{Independent random errors} reduce in the combined estimate
when both methods are averaged; the agreement between them tightens
as both individual uncertainties tighten. \engterm{Independent
systematic errors} cancel partially in the difference: a calliper
that reads $20 \,\si{\micro\meter}$ high and a laser that reads $20
\,\si{\micro\meter}$ low produce a difference of $40 \,\si{\micro
\meter}$, which the compatibility ratio flags. \engterm{Common
systematic errors} go undetected: a temperature bias affecting both
methods in the same direction shifts both readings together, and the
difference is unchanged. Cross-checking detects independent errors;
common errors require an independent control on the shared cause.

\subsection{Worked example: the volume of a stone}

A reader measures the volume of an irregular stone by water
displacement and by geometric approximation. The displacement gives
$V_{A} = 47.3 \,\si{\milli\liter}$ with $u_{A} = 1.0 \,\si{\milli
\liter}$ derived from the graduated cylinder's resolution and the
parallax of the meniscus reading. The geometric approximation,
modelling the stone as a triaxial ellipsoid with semi-axes
$2.6 \pm 0.1$, $2.0 \pm 0.1$, and $1.8 \pm 0.1 \,\si{\centi\meter}$,
gives $V_{B} = \tfrac{4}{3}\pi(2.6)(2.0)(1.8) \approx 39.2 \,\si
{\milli\liter}$. The product-of-powers rule gives the relative
uncertainty of $V_{B}$ as $\sqrt{(0.1/2.6)^{2} + (0.1/2.0)^{2} +
(0.1/1.8)^{2}} \approx \sqrt{0.0015 + 0.0025 + 0.0031} \approx 0.084$
or $8.4\,\%$, so $u_{B} \approx 3.3 \,\si{\milli\liter}$.

The compatibility ratio:

\[
r = \frac{|47.3 - 39.2|}{\sqrt{1.0^{2} + 3.3^{2}}} =
\frac{8.1}{3.4} \approx 2.4.
\]

The two methods disagree at the $2.4\sigma$ level. The disagreement
is large enough to flag a previously uncharacterised error source.
Two candidates are immediate. The geometric model may be too smooth,
because the stone has ridges and concavities the ellipsoid does not
capture. The displacement reading may also be biased high if air
bubbles clung to the stone's surface. The next investigation step
is to wet the stone thoroughly
before displacement, repeat both measurements, and either close the
gap or find a third independent method to break the tie.

The third method that closes ties in this chapter's project is
image-based or contact-based dimensional reconstruction. We will
return to this worked example in the project specification below.

\subsection{The discipline, restated}

\begin{principle}[Cross-checking is the wrong-ruler detector]
A single measurement is consistent with itself by construction. Two
independent measurements either agree within their combined
uncertainty or disagree. Disagreement is information: it locates an
error source that a single measurement's uncertainty budget did not
include. The discipline of cross-checking is the only detector of a
wrong ruler that does not depend on knowing the right answer in
advance.
\end{principle}

The principle has a corollary. A measurement reported alone, without
an independent check, has only its self-derived uncertainty to defend
it. The reader who takes a single number as authoritative because it
came from a precise instrument has confused precision with accuracy.
Hubble's primary mirror was figured to high precision against a
test instrument that was itself wrong; the precision of the figure
was not the issue. The wrong-ruler detector that should have caught
the error was the auxiliary tests; they had been built into the
process and were overruled. We turn to that case now.

\section{Failure: when the ruler was wrong (Hubble's primary
mirror)}\pathtag{core}

A length measurement chain is only as honest as its least independent
test. The Hubble Space Telescope primary mirror was figured between
1979 and 1981 against a reflective null corrector that was itself
mispositioned, and the mirror inherited the error. Auxiliary tests
indicated the discrepancy at fabrication time and were set aside.
The case has been treated in chapter 2 of this volume as a
metrology-chain failure, in chapter 3 as a test-instrument-as-
failed-standard case, and in chapter 4 as an auxiliary-test override
case. We return here to the narrowest framing this chapter can use:
the ruler against which the mirror was measured was wrong, and no
internal check inside the ruler's measurement chain could have
found the error.

\subsection{The mechanism}

The Hubble Space Telescope launched aboard Space Shuttle Discovery
on 24 April 1990. Imaging anomalies were identified during the first
weeks of on-orbit operation. The NASA Optical Systems Failure
Report, prepared by the Optical Systems Board of Investigation
chaired by Lew Allen Jr.\ and released in November 1990, identified
the cause as a misalignment in the reflective null corrector (RNC),
the optical-test instrument used to measure the mirror's figure
during polishing. A field lens within the RNC was mispositioned along
the optical axis by approximately $1.3 \,\si{\milli\meter}$. The
mispositioning resulted from a misinterpretation of a measurement
target on a metering rod used during RNC assembly: the end cap and
the target feature were not properly distinguished, and the spacing
was set against the wrong feature. The error in the test instrument
propagated into the polished mirror's figure: the mirror's figure
error corresponded to severe spherical aberration, described by the
investigation report as a $0.4$-wave rms wavefront error at $632.8
\,\si{\nano\meter}$; in the local short form, the outer edge was
about $2 \,\si{\micro\meter}$ too flat, an error large enough to
ruin diffraction-limited imaging across the optical bandpass
\cite{acc:nasa-hubble-optical-systems-1990}.

The case is recorded in the named-cases registry at
\repopath{docs/research/accidents/hubble-primary-mirror-1990.md}.

\begin{figure}[ht]
\centering
\begin{tikzpicture}[x=1cm, y=1cm, font=\small]

\node[anchor=south, font=\bfseries] at (7, 6.2) {Hubble primary mirror: the wrong-ruler chain};

% == The internal chain (top row)
\node[draw, rounded corners, thick, minimum width=2.5cm, minimum height=0.9cm, fill=engblue!12]
  (rod) at (1.6, 4.6) {metering rod};
\node[draw, rounded corners, thick, minimum width=2.5cm, minimum height=0.9cm, fill=engblue!12]
  (rnc) at (5.4, 4.6) {RNC assembly};
\node[draw, rounded corners, thick, minimum width=2.5cm, minimum height=0.9cm, fill=engblue!12]
  (figref) at (9.2, 4.6) {RNC reference figure};
\node[draw, rounded corners, thick, minimum width=2.5cm, minimum height=0.9cm, fill=engblue!12]
  (mirror) at (13.0, 4.6) {polished mirror};

\draw[->, thick] (rod) -- (rnc) node[midway, above, font=\scriptsize] {sets};
\draw[->, thick] (rnc) -- (figref) node[midway, above, font=\scriptsize] {carries};
\draw[->, thick] (figref) -- (mirror) node[midway, above, font=\scriptsize] {polished to};

% Error injected at the metering-rod target
\node[anchor=south, font=\scriptsize, engred] at (1.6, 5.6) {target misread};
\node[anchor=south, font=\scriptsize, engred] at (1.6, 5.95) {by $1.3 \,\si{\milli\meter}$};
\draw[->, thick, engred] (1.6, 5.5) -- (1.6, 5.1);

% Error propagated to the mirror
\node[anchor=north, font=\scriptsize, engred] at (13.0, 4.1) {outer edge};
\node[anchor=north, font=\scriptsize, engred] at (13.0, 3.7) {$\approx 2 \,\si{\micro\meter}$ too flat};

% == Internal consistency check labels (between top boxes)
\node[anchor=north, font=\scriptsize, gray, align=center] at (3.5, 4.0) {internal\\check $\checkmark$};
\node[anchor=north, font=\scriptsize, gray, align=center] at (7.3, 4.0) {internal\\check $\checkmark$};
\node[anchor=north, font=\scriptsize, gray, align=center] at (11.1, 4.0) {internal\\check $\checkmark$};

% == Auxiliary tests (bottom)
\node[draw, rounded corners, thick, minimum width=3.2cm, minimum height=0.9cm, fill=engred!12]
  (refractive) at (5.4, 1.4) {refractive null corrector};
\node[draw, rounded corners, thick, minimum width=3.2cm, minimum height=0.9cm, fill=engred!12]
  (inverse) at (9.2, 1.4) {inverse null test};

% Disagreement arrows (auxiliary tests vs figref)
\draw[->, thick, engred, densely dashed] (refractive.north) -- ([yshift=-0.0cm]figref.south west) node[midway, left, font=\scriptsize] {disagrees};
\draw[->, thick, engred, densely dashed] (inverse.north) -- ([yshift=-0.0cm]figref.south east) node[midway, right, font=\scriptsize] {disagrees};

% "Overruled" annotation
\node[draw, dashed, thick, engred, font=\scriptsize, fill=white] at (7.3, 0.4) {disagreements set aside under cost / schedule pressure};

\end{tikzpicture}
\caption[The Hubble primary mirror as a wrong-ruler
chain]{The Hubble primary mirror's measurement chain illustrates the
wrong-ruler pattern. Each internal step in the polishing chain
(metering rod, RNC assembly, RNC reference figure, polished mirror)
checked the link above it and reported agreement. The error in the
chain was injected at the top, when a measurement target on the
metering rod was misinterpreted by about $1.3 \,\si{\milli\meter}$,
and the entire chain inherited the error. Two auxiliary optical tests
(the refractive null corrector and the inverse null test) disagreed
with the RNC's reference figure at fabrication time; the
disagreements were set aside, and the wrong-ruler detector that
might have caught the error was disabled. Source: editors' schematic
following the OSBI report
\cite{acc:nasa-hubble-optical-systems-1990}.}
\label{fig:vol01:ch07:hubble-wrong-ruler}
\end{figure}


\Cref{fig:vol01:ch07:hubble-wrong-ruler} maps the chain. The Hubble
mirror appears in chapters 2--4 of this volume as a metrology-chain
case, a test-instrument-as-failed-standard case, and an
auxiliary-test override case. The framing this chapter adds is
specific to length measurement: the mirror's figure was a length
measurement at the optical-wavelength scale, the RNC was the length
reference against which it was made, and the chain's blind spot was
the assumption that the reference's own length could not be checked
against another length instrument at the same scale.

\paragraph{Technical mechanism.}
The polished mirror's figure was measured against the RNC's
reference figure to within the polishing process's uncertainty. The
RNC's reference figure was produced by the geometric placement of
its optical elements, in particular a field lens whose axial
position was set against a measurement target on a metering rod.
The metering rod had two features near its end that the assembler
could mistake for one another: an end cap and a target feature, the
latter recessed by approximately $1.3 \,\si{\milli\meter}$ from the
former. The lens was set against the end cap rather than the target
feature; the RNC was thereafter internally consistent with the
wrong axial spacing; the mirror was polished to match the RNC's
wrong figure. The error in length at the metering-rod input scaled
through the optical-test geometry into a $0.4$-wave rms wavefront
error at $632.8 \,\si{\nano\meter}$, or about $2 \,\si{\micro\meter}$
of figure error at the mirror's outer edge
\cite{acc:nasa-hubble-optical-systems-1990}.

\paragraph{Organisational context.}
The Allen Commission's analysis identified the procedural barrier
that should have caught the error. Two auxiliary optical tests, a
refractive null corrector and an inverse null test, both indicated
discrepancies with the RNC's reading at fabrication time. The
process treated the RNC as the more authoritative instrument, on
the basis of its design pedigree and a quality-assurance regime
that recorded auxiliary-test anomalies but did not require their
resolution. Cost and schedule pressure during the final polishing
phase reinforced the disposition to set the auxiliary results
aside. The Commission's recommendations centred on the missing
procedural barriers: independent verification of the test
instrument, anomaly-resolution procedures that cannot be informally
overridden, and engaged customer (NASA) oversight of the
optical-test apparatus rather than only of the test results
\cite{acc:nasa-hubble-optical-systems-1990}.

\paragraph{Lesson for length measurement.}
The RNC was a length-measurement instrument whose own length-chain
provenance was not independently audited. For a working length
measurement, the lesson takes the operational form developed in
section 7.5: a reference instrument is itself a measurement chain,
and the cross-check on the reference is an independent measurement
of the reference, not a repeat of the same measurement on the same
chain. A calliper whose calibration is verified against a gauge
block calibrated against the same master, returned from the same
calibration laboratory, is not cross-checked; it is internally
consistent. The independent check is a measurement against a
different master, from a different chain, with a different
calibration history. The discipline is no more elaborate than
that; what makes it hard is procedural, not technical, exactly as
in the Hubble case.

\subsection{The wrong-ruler diagnosis}

The Hubble mirror's measurement chain looks straightforward. The
mirror was polished against a reference figure; the reference figure
was carried by the reflective null corrector; the RNC was assembled
according to a specification that placed each optical element at a
known axial position, set against a metering rod whose features
defined the spacings. Every link in the chain produced an internally
consistent measurement. The mirror's figure agreed with the RNC's
reference figure to within the polishing process's uncertainty. The
RNC's assembly agreed with its specification to within the metering
rod's uncertainty.

What no internal check could detect was that the metering-rod
reading was misinterpreted. The end cap and the target feature were
both legitimate features on the rod; the assembler set the field
lens against the wrong one. The downstream chain, asked only whether
each link agreed with the link above it, had no way to flag that the
top of the chain was set against the wrong reference.

The chain that could have caught the error was the auxiliary
tests: a refractive null corrector and an inverse null test, both of
which had indicated the discrepancy at fabrication time. The
refractive test produced a result inconsistent with the RNC's
reading; the inverse test produced a similarly inconsistent
indication. The fabrication process treated the RNC as the
authoritative instrument and set the auxiliary results aside under
cost and schedule pressure \cite{acc:nasa-hubble-optical-systems-1990}.

The wrong-ruler diagnosis from this chapter's vocabulary is direct.
The RNC was a length measurement chain whose internal checks were
self-consistent and whose external checks (the auxiliary tests)
disagreed. A compatibility ratio computed at fabrication time, with
the auxiliary tests as independent methods, would have flagged the
RNC's reading as incompatible with the auxiliary results. The
process did not run that comparison, or ran it and overruled the
result.

\subsection{What the case tells us}

The Hubble mirror is a case in which precision was very high and
accuracy was very low. The polishing process was capable of figuring
the mirror to a small fraction of a wavelength of visible light;
that capability was used to figure the mirror to the wrong shape.
Precision did not save the measurement, because the measurement's
reference was wrong. A reader who treats a precise instrument as
authoritative without independent verification repeats the Hubble
failure mode.

The recovery, COSTAR (Corrective Optics Space Telescope Axial
Replacement), installed during the STS-61 servicing mission in
December 1993, applied a corrective optic that compensated for the
mirror's known figure error. The correction worked because, after
launch, the figure error had been characterised independently using
the on-orbit imaging itself; once the error was known, a corrective
optic could be designed to undo it. The pre-launch failure was not
that the figure could not be corrected; it was that the figure
error was not known before launch.

\subsection{The discipline, applied}

\begin{principle}[A reference instrument is itself a measurement chain]
Any instrument used to measure another instrument is itself a
measurement chain, with its own uncertainty, its own calibration
history, and its own opportunities for error. A reference instrument
is verified independently of the measurement it is used to perform;
when independent verification disagrees with the reference's
indication, the disagreement is investigated, not overruled. A
process that lets an authoritative instrument's reading silence
disagreeing auxiliary tests has built a wrong-ruler detector and
disabled it.
\end{principle}

The principle restates, at the metrology-chain scale, the cross-
checking principle of section 7.5. The two are the same principle
applied at different layers: the auxiliary tests are the cross-check
on the reference instrument, just as the second method is the
cross-check on the first. Both fail in the same way, when the
disagreement is treated as an inconvenience rather than as
information.

\section{Project}\pathtag{core}

\begin{project}[The volume of an irregular object, by three methods]
The reader picks one irregular household object and measures its
volume by three independent methods. Each measurement carries an
uncertainty bar derived from the chapter 4 algebraic rules. The
three values are compared. Where the methods agree within their
combined uncertainty, the reader's confidence in the volume grows;
where they disagree, the reader investigates which method's
uncertainty budget is incomplete.

\textbf{Track}. Household and analysis (hybrid). The reader works at
the kitchen table or workbench; no climbing, no chemical work, no
electrical work above a battery. \textbf{Hazard class}: none,
provided the reader chooses an expendable object, uses stable
glassware or plasticware, and keeps water away from electrical
devices. The chosen object should not be sharp, electrical,
valuable, or water-damaged. \textbf{Tools}. A kitchen scale or
postal scale, resolution about $1 \,\si{\gram}$. A ruler or tape
measure, resolution about $1 \,\si{\milli\meter}$. A beaker,
measuring cup, or graduated cylinder for water displacement. Tap
water. A smartphone camera, optional for image-based
reconstruction. Paper and pencil for the geometric approximation.

\textbf{Choosing the object}. Pick one irregular object whose volume
is roughly between $5$ and $200 \,\si{\milli\liter}$ and which is
not soluble, absorbent, fragile, or buoyant in water. Suggested
objects: a smooth stone from a garden or a beach, a small piece of
fruit (an apple or a small pear, intact), a small hand tool (a
spanner, a pair of pliers), a piece of jewellery the reader is
willing to immerse, a piece of solid pottery, a candle. Avoid sponge,
foam, soft fruit with cut surfaces, paper, or porous wood; these
either soak up water or float.

\textbf{Procedure}. The reader performs three of the four methods
below. Method (i), water displacement, is required. The reader
chooses two of the remaining three.

\begin{enumerate}[leftmargin=*]
  \item \engterm{Water displacement (Archimedes)}. Required.
    Fill the graduated cylinder or measuring cup with water to a
    measured initial level $V_{0}$. Submerge the object fully,
    ensuring no air bubbles cling to its surface (gentle agitation
    or pre-wetting helps). Read the new water level $V_{1}$. The
    object's volume is $V = V_{1} - V_{0}$. Repeat the measurement
    five times, removing and re-immersing the object between
    readings. Record the five readings; report $V$ as the mean and
    $u(V)$ as the standard deviation of the five readings combined
    in quadrature with the graduated cylinder's resolution-derived
    uncertainty (use a uniform-distribution standard uncertainty
    of half the graduation interval divided by $\sqrt{3}$ per
    chapter 4).
  \item \engterm{Geometric approximation}. Decompose the object into
    a sum (and difference) of simple shapes whose volumes the
    reader can compute from length measurements: spheres, cylinders,
    cones, rectangular boxes, ellipsoids. For each component shape,
    measure its defining lengths to the resolution of the ruler or
    calliper available; record each length with its standard
    uncertainty (resolution divided by $\sqrt{3}$ for a uniform
    distribution, plus any repeatability spread). Compute each
    component volume and propagate its uncertainty by the chapter 4
    product-of-powers rule. Sum the component volumes (subtracting
    any voids); the standard uncertainty of the sum is the
    quadrature sum of the component standard uncertainties when the
    components are uncorrelated. Report $V$ and $u(V)$.
  \item \engterm{Image-based reconstruction}. With the smartphone
    camera, take photographs of the object from at least eight
    angles around its perimeter and from at least two heights, with
    a known reference object (a coin or a ruler) in each frame for
    scale. Use a free or commercial photogrammetry tool, or a
    smartphone scanning app, to construct a three-dimensional
    surface mesh and report its volume. Estimate the method
    uncertainty by repeating the reconstruction at least three
    times: once with the nominal scale reference, once with the
    scale reference shifted by its estimated uncertainty, and once
    with a different mesh-resolution setting. Report the spread as
    a sensitivity check, and state plainly that this is a method
    repeatability estimate rather than a traceable photogrammetry
    uncertainty.
  \item \engterm{Contact-based dimensional reconstruction}.
    \textit{Mastery option.} Sample the object's surface at a
    regular grid of points using a calliper or ruler. Record enough
    cross-sections to approximate the object as stacked slices,
    compute the area of each slice, and estimate volume by summing
    slice area times slice spacing. Repeat the sampling; the spread
    is a repeatability estimate. A closed-mesh divergence-theorem
    calculation belongs in Volume II or in a mastery note.
\end{enumerate}

\textbf{Record}. Each method's record contains: the chosen object;
the date of the measurements; the instrument(s) used and their
resolution; the conditions (water temperature, ambient temperature,
the number of repeats); the raw readings; the computed value $V$;
the propagated standard uncertainty $u(V)$; the dominant uncertainty
source identified by sensitivity analysis. Where the method involves
software, record the software, version, settings, and the algorithm
used to compute the volume.

\textbf{Analysis}. Tabulate the three methods' results as a small
table: method, value, expanded uncertainty at $k = 2$, dominant
uncertainty source. Compute each pairwise compatibility ratio
$r_{AB}, r_{AC}, r_{BC}$ using the formula of section 7.5. Identify
which pairs agree within $r \leq 2$ and which do not. For pairs that
disagree, propose at least one previously uncharacterised error
source that could explain the disagreement, and propose a
diagnostic measurement that would either confirm or refute the
proposed source.

\textbf{Reflection (1500 words)}. The reader writes a 1500-word
narrative addressing:

\begin{itemize}[leftmargin=*]
  \item Which method gave the smallest expanded uncertainty, and
    which input dominated that uncertainty.
  \item Where two methods agreed within their combined uncertainty,
    and where they did not. For pairs that disagreed, what
    additional measurement would the reader propose to investigate
    the discrepancy.
  \item Which method's uncertainty budget the reader is least
    confident about, and what would be needed to tighten it.
  \item One household object whose volume the reader's chosen
    methods cannot adequately measure, and what method would be
    required.
  \item One measurement in the reader's life or work where the
    same three-method discipline would change the reader's
    confidence in the result.
\end{itemize}

\textbf{Deliverable}. The three method records, the comparison
table, the compatibility ratios, and the reflection. The general
editor's reference set of three-method computations for a
representative stone, with worked uncertainty arithmetic, is in
\repopath{volumes/01-quantity/ch07-length-area-volume-mass/data/stone\_volume\_three\_methods.csv};
the reader is encouraged to consult it \emph{after} completing their
own measurements rather than before.

\textbf{Time}. One to two hours per method, plus two to three hours
of analysis and writing. Total approximately five to nine hours,
distributed across one or two days.

\subsubsection*{Phases, deliverables, and success criteria}

The project decomposes into five phases, each with a concrete
deliverable and a success criterion the reader can self-check
before moving on.

\begin{enumerate}[leftmargin=*]
  \item \engterm{Phase 1: planning} (about 30 minutes).
    \textit{Deliverable}: a one-page measurement plan that names the
    object, the three methods chosen, the instruments and their
    stated resolutions, the expected dominant uncertainty source per
    method, and a target expanded uncertainty at $k = 2$ for each
    method.
    \textit{Success criterion}: a peer reading the plan can predict
    the order of magnitude of each method's expanded uncertainty.
  \item \engterm{Phase 2: data collection} (about 1--2 hours per
    method). \textit{Deliverable}: for each method, the raw readings
    (at least five per method), the recorded environmental
    conditions (water temperature, room temperature, instrument
    temperature where relevant), and the timestamp of each session.
    \textit{Success criterion}: a reader returning to the data three
    months later can reproduce the analysis from the raw record
    without consulting memory.
  \item \engterm{Phase 3: uncertainty propagation} (about 1 hour).
    \textit{Deliverable}: for each method, a written uncertainty
    budget table with rows for each input quantity, the
    distributional assumption (Gaussian, uniform, triangular), the
    standard uncertainty, and the contribution to the combined
    standard uncertainty.
    \textit{Success criterion}: removing the dominant input from the
    table and recomputing shows the combined standard uncertainty
    drop noticeably; removing a sub-dominant input changes the
    combined standard uncertainty by a factor less than $1.1$.
  \item \engterm{Phase 4: cross-checking} (about 30 minutes).
    \textit{Deliverable}: a three-row table (method, $V \pm U_{k=2}$,
    dominant source), the three pairwise compatibility ratios
    $r_{AB}, r_{AC}, r_{BC}$ per section 7.5, and a short narrative
    naming pairs that agree and pairs that disagree.
    \textit{Success criterion}: each $r$ value is computed correctly
    from the table; for any pair with $r > 2$, a proposed mechanism
    and a proposed diagnostic measurement are documented.
  \item \engterm{Phase 5: reflection} (about 1--2 hours).
    \textit{Deliverable}: the 1500-word reflection specified above.
    \textit{Success criterion}: the reflection answers each of the
    five prompted questions with a specific measurement, not a
    generic discussion of measurement principles.
\end{enumerate}

\subsubsection*{Worked example: a stone, three methods}

A reference set of three measurements for a representative stone,
performed by the editors, is recorded in
\repopath{volumes/01-quantity/ch07-length-area-volume-mass/data/stone\_volume\_three\_methods.csv}.
The displacement method gave $V_{A} = 47.3 \,\si{\milli\liter}$
with $U_{A} = 2.0 \,\si{\milli\liter}$ at $k = 2$, dominated by
parallax and resolution. The geometric ellipsoid approximation gave
$V_{B} = 39.2 \,\si{\milli\liter}$ with $U_{B} = 6.6 \,\si{\milli
\liter}$ at $k = 2$, dominated by the model's fit to the actual
ridged shape. The smartphone photogrammetry gave $V_{C} = 44.8 \,
\si{\milli\liter}$ with $U_{C} = 3.4 \,\si{\milli\liter}$ at $k = 2$,
dominated by the photogrammetric scale-reference uncertainty. The
pairwise compatibility ratios are $r_{AB} \approx 2.4$ (disagree),
$r_{AC} \approx 1.8$ (borderline), $r_{BC} \approx 1.5$ (agree).
The geometric method's disagreement with the displacement method is
the most informative result: it locates a previously
uncharacterised model-fit error in the ellipsoid approximation
larger than the uncertainty budget had carried. The reader's own
measurements will produce a different pattern; the discipline that
makes the project a measurement experience and not a sequence of
arithmetic exercises is the comparison of \emph{their} numbers,
honestly recorded.
\end{project}

\section{Exercises}

The exercises train calculation, derivation, estimation, simulation,
design, diagnosis, failure analysis, and judgment. Calculation,
derivation, and computational exercises have full solution sketches;
design, judgment, and other open-ended exercises have discussion
rubrics that name the points a defensible reply must cover.
Exercises are tagged with their reader-path tier.

\subsection*{Calculation}\pathtag{core}

\begin{exercise}
A cylinder has radius $R = 25.0 \,\si{\milli\meter} \pm 0.1 \,\si
{\milli\meter}$ and height $H = 80.0 \,\si{\milli\meter} \pm 0.5 \,
\si{\milli\meter}$. Compute the volume $V = \pi R^{2} H$ and its
relative standard uncertainty. Identify which input contributes more
to $u(V)/V$ and explain why the radius's exponent doubles its
contribution before quadrature.
\end{exercise}

\paragraph{Solution.}
The calculation proceeds in six steps that the reader carries
through any product-of-powers uncertainty problem.

\emph{Step 1: compute the nominal value.}
$V = \pi R^{2} H = \pi \times (25.0)^{2} \times 80.0 = \pi \times
625 \times 80 = 50\,000\pi \approx 157\,080 \,\si{\milli\meter
\cubed} = 157.1 \,\si{\milli\liter}$.

\emph{Step 2: compute the relative input uncertainties.}
$u(R)/R = 0.1 / 25.0 = 0.0040 = 0.40\,\%$;
$u(H)/H = 0.5 / 80.0 = 0.00625 = 0.625\,\%$.

\emph{Step 3: apply the product-of-powers rule with the
exponents on $R$ and $H$ in $V = \pi R^{2} H$.}
The exponents are $+2$ on $R$ and $+1$ on $H$, so
\[
\left(\frac{u(V)}{V}\right)^{2}
  = 2^{2} \left(\frac{u(R)}{R}\right)^{2}
  + 1^{2} \left(\frac{u(H)}{H}\right)^{2}
  = 4 \, (0.0040)^{2} + (0.00625)^{2}.
\]
The factor $\pi$ is exact and contributes no uncertainty, as does
the integer $2$ in $R^{2}$.

\emph{Step 4: evaluate the variance terms separately, so the
contributions are visible.}
Radius term: $4 \times 1.6 \times 10^{-5} = 6.4 \times 10^{-5}$.
Height term: $3.9 \times 10^{-5}$. Sum: $1.03 \times 10^{-4}$.
The radius term contributes $62\,\%$ of the variance, the height
term $38\,\%$.

\emph{Step 5: take the square root and convert to absolute
uncertainty.}
$u(V)/V = \sqrt{1.03 \times 10^{-4}} \approx 0.01016$, so
$u(V) \approx 0.0102 \times 157.1 \approx 1.60 \,\si{\milli
\liter}$.

\emph{Step 6: report with the right number of digits, and explain
the radius's dominance.}
Report $V = 157.1 \pm 1.6 \,\si{\milli\liter}$ at $k = 1$, or
$V = 157 \pm 2 \,\si{\milli\liter}$ rounded to the working
precision. The radius dominates because the product-of-powers rule
weights each term by the squared exponent before quadrature, and
the squared exponent on $R$ is $4$ whereas on $H$ it is $1$. A
$1\,\%$ relative uncertainty on $R$ enters the variance as
$4 \times 10^{-4}$; the same $1\,\%$ on $H$ enters as $1 \times
10^{-4}$. The lesson generalises: in a product-of-powers formula,
the variable with the largest absolute exponent is the variable
whose precision deserves the largest investment, even when its
relative uncertainty is the smaller of the two.

\begin{exercise}
An object has measured mass $m = 138.5 \,\si{\gram} \pm 0.5 \,\si
{\gram}$ and measured volume $V = 17.4 \,\si{\milli\liter} \pm
0.6 \,\si{\milli\liter}$, both at $k = 1$. Compute the density and
its relative standard uncertainty. Identify which of (a) doubling
the precision of the mass measurement, (b) doubling the precision of
the volume measurement, would more reduce the density's
uncertainty.
\end{exercise}

\paragraph{Solution.}

\emph{Step 1: nominal density.}
$\rho = m / V = 138.5 / 17.4 = 7.96 \,\si{\gram\per\centi\meter
\cubed}$. The published density of plain carbon steel is
$7.85 \,\si{\gram\per\centi\meter\cubed}$
\cite{text:ashby2017}; the measured value is consistent with steel
within the uncertainty we are about to compute, and inconsistent
with aluminium ($2.70$), copper ($8.96$), and lead ($11.34$) by
many standard deviations.

\emph{Step 2: relative input uncertainties.}
$u(m)/m = 0.5 / 138.5 = 0.0036 = 0.36\,\%$;
$u(V)/V = 0.6 / 17.4 = 0.0345 = 3.45\,\%$.

\emph{Step 3: product-of-powers rule with exponents $+1$ and $-1$.}
For $\rho = m \cdot V^{-1}$ with uncorrelated $m$ and $V$,
\[
\left(\frac{u(\rho)}{\rho}\right)^{2} = (+1)^{2}
\left(\frac{u(m)}{m}\right)^{2} + (-1)^{2}
\left(\frac{u(V)}{V}\right)^{2}.
\]
The sign of the exponent vanishes in the squared sum; the rule
sums variances regardless of whether the variable is in the
numerator or the denominator.

\emph{Step 4: evaluate each variance term.}
Mass term: $0.0036^{2} = 1.30 \times 10^{-5}$.
Volume term: $0.0345^{2} = 1.19 \times 10^{-3}$.
Sum: $1.20 \times 10^{-3}$. The volume contributes
$99\,\%$ of the variance; the mass contributes $1\,\%$. The
disparity is the answer to the question and motivates the next
step.

\emph{Step 5: absolute uncertainty on density.}
$u(\rho)/\rho = \sqrt{1.20 \times 10^{-3}} \approx 0.0346$, so
$u(\rho) \approx 0.0346 \times 7.96 \approx 0.28 \,\si{\gram\per
\centi\meter\cubed}$. Report $\rho = 7.96 \pm 0.28 \,\si{\gram\per
\centi\meter\cubed}$ at $k = 1$, or $\rho = (8.0 \pm 0.3) \,\si
{\gram\per\centi\meter\cubed}$.

\emph{Step 6: the precision-budget question.}
(a) Doubling the precision of the mass measurement halves $u(m)$,
so $u(m)/m$ falls from $0.0036$ to $0.0018$. The mass variance
term falls from $1.30 \times 10^{-5}$ to $3.2 \times 10^{-6}$. The
new combined variance is $1.19 \times 10^{-3} + 3.2 \times 10^{-6}
\approx 1.19 \times 10^{-3}$: indistinguishable from the original
to three significant digits. The improvement on $u(\rho)/\rho$
is at the fifth decimal place.
(b) Doubling the precision of the volume measurement halves
$u(V)$, so $u(V)/V$ falls from $0.0345$ to $0.0173$. The volume
variance term falls from $1.19 \times 10^{-3}$ to $2.99 \times
10^{-4}$. The combined variance is $2.99 \times 10^{-4} + 1.30
\times 10^{-5} \approx 3.12 \times 10^{-4}$, so $u(\rho)/\rho
\approx 0.0177$ and $u(\rho) \approx 0.14 \,\si{\gram\per\centi
\meter\cubed}$. The density's relative uncertainty halves.
The general rule reads off this calculation: spend precision
budget where the largest relative uncertainty lives, because
quadrature combination ignores small terms whose squared
contribution is far below the dominant term.
\repopath{volumes/01-quantity/ch07-length-area-volume-mass/code/density\_propagation.py}
encodes this and also checks alloy compatibility.

\begin{exercise}
A length measured with a calliper at $20.0 \,\si{\degreeCelsius}$
reads $1.0000 \,\si{\meter}$. The workpiece is steel with thermal
expansion coefficient $\alpha = 1.2 \times 10^{-5} \,\si{\per
\kelvin}$. Compute the corrected length at a workshop temperature of
$25.0 \,\si{\degreeCelsius}$, with the assumption that the calliper
itself does not change length. State the bias the uncorrected
reading introduces and identify the engineering decision in which
this bias is non-negligible.
\end{exercise}

\paragraph{Solution.}

\emph{Step 1: write down the linear-expansion model.}
A material at reference temperature $T_{0}$ with reference length
$L_{0}$ has length $L(T) = L_{0} [1 + \alpha (T - T_{0})]$ at
temperature $T$, with $\alpha$ the coefficient of linear thermal
expansion. For small fractional changes the model linearises to
$\Delta L = L_{0} \alpha \Delta T$, the working form used here.

\emph{Step 2: substitute the workpiece values.}
$L_{0} = 1.0000 \,\si{\meter}$; $\alpha_{\mathrm{steel}} = 1.2
\times 10^{-5} \,\si{\per\kelvin}$ (a representative carbon-steel
value; alloy steels run $1.0$ to $1.7 \times 10^{-5} \,\si{\per
\kelvin}$, and the working tolerance budget should carry the
material-specific value when it is known); $\Delta T = 5.0 \,\si
{\kelvin}$.
\[
\Delta L = 1.0000 \times 1.2 \times 10^{-5} \times 5.0 = 6.0
\times 10^{-5} \,\si{\meter} = 60 \,\si{\micro\meter}.
\]

\emph{Step 3: state the corrected length and the calliper
assumption.}
The corrected length at the workshop temperature is $L(25) =
L_{0} + \Delta L = 1.00006 \,\si{\meter}$, provided the calliper
itself is at the reference temperature $T_{0} = 20.0 \,\si
{\degreeCelsius}$. If the calliper has also warmed to $25.0 \,\si
{\degreeCelsius}$ and the calliper material is steel of the same
$\alpha$, the calliper has expanded by the same fractional amount,
and the calliper reading at the warm temperature reports the
corrected length directly. The classic mixed case is a steel
workpiece measured with a calliper made of a different metal:
hardened tool steels have $\alpha$ near $1.0 \times 10^{-5}$, and
the calliper-and-workpiece mismatch produces a residual bias
proportional to the difference in $\alpha$.

\emph{Step 4: identify the regime where the bias matters.}
A $60 \,\si{\micro\meter}$ bias on a metre is $60$ parts per
million, or $0.006\,\%$. The bias matters in absolute terms
whenever the dimensional tolerance approaches the bias's
magnitude. The ISO 286 system of fits records canonical tolerance
grades for shaft-and-bore fits: at a $1 \,\si{\meter}$ size, the
IT7 grade is $\pm 40 \,\si{\micro\meter}$, IT6 is $\pm 22 \,\si
{\micro\meter}$, and IT5 is $\pm 16 \,\si{\micro\meter}$. The
worked $\Delta T = 5 \,\si{\kelvin}$ bias of $60 \,\si{\micro
\meter}$ already exceeds the IT7 tolerance band on its own, and
multiples of it exceed IT5 and IT6 by several-fold.

\emph{Step 5: name the engineering decision in which the bias is
non-negligible.}
Any fit of class H7/g6 or tighter on a precision shaft of order
$1 \,\si{\meter}$ length cannot be measured without thermal
correction in a workshop above the standard reference temperature.
Gauge-block calibration laboratories run at $20.0 \,\si{\degree
Celsius} \pm 0.5 \,\si{\kelvin}$ for exactly this reason
\cite{text:doiron-beers-2009}. A machinist who reports a length
without recording the temperature of measurement is reporting a
length whose uncertainty is unknowable, because the dominant
source of bias is not in the instrument reading but in the
unrecorded $\Delta T$.

\emph{Step 6: invoke the standard.}
ISO 1 mandates correction to $20 \,\si{\degreeCelsius}$ for
dimensional specifications precisely to remove this bias from the
documented length \cite{std:iso-1-2022}. The standard does not
require that measurements be made at $20 \,\si{\degreeCelsius}$;
it requires that lengths reported per the specification be the
lengths the part would have at $20 \,\si{\degreeCelsius}$. A
length reported per ISO 1 at any other temperature is a length
whose correction back to $20$ has been applied and recorded.

\begin{exercise}
A water-displacement measurement reports a volume change of
$V = 47.5 \,\si{\milli\liter}$ on a graduated cylinder whose
graduations are at $1 \,\si{\milli\liter}$ intervals. Treat the
reading as uniformly distributed within $\pm 0.5 \,\si{\milli\liter}$
of the indicated graduation and compute the standard uncertainty
arising from resolution alone. Then add an independent reading-spread
component of $\pm 0.3 \,\si{\milli\liter}$ standard deviation
estimated from five repeats and report the combined standard
uncertainty.
\end{exercise}

\paragraph{Solution.}
The water-displacement measurement combines a Type B resolution
contribution with a Type A spread contribution. The solution
proceeds through the resolution calculation, the two-reading
combination, the addition of the spread, and the final report.

\emph{Step 1: convert the resolution to a standard uncertainty.}
The graduated cylinder reads to $1\,\si{\milli\liter}$ intervals.
The exercise treats the reading as uniformly distributed within
$\pm 0.5\,\si{\milli\liter}$ of the indicated graduation. The
standard deviation of a uniform distribution on $[-a, +a]$ is
$a/\sqrt{3}$, so the single-reading resolution-derived standard
uncertainty is
\[
u_{\text{res, single}} = \frac{0.5}{\sqrt{3}} \approx 0.289\,\si{
\milli\liter}.
\]

\emph{Step 2: account for the two independent readings in the
displacement.} The volume change is the difference of two
graduated-cylinder readings (before and after immersion). Each
reading carries its own independent resolution uncertainty. The
two-reading difference has variance equal to the sum of the
individual variances:
\[
u_{\text{res, diff}}^{2} = u_{\text{res, single}}^{2} +
u_{\text{res, single}}^{2} = 2 u_{\text{res, single}}^{2},
\]
so
\[
u_{\text{res, diff}} = \sqrt{2} \times 0.289 \approx 0.408\,\si{
\milli\liter}.
\]
The $\sqrt{2}$ factor is the price of a difference measurement: the
volume is computed from two readings, and the noise from both
contributes.

\emph{Step 3: characterise the Type A spread.} The problem states a
$0.3\,\si{\milli\liter}$ single-reading spread from five repeats.
This is the population standard deviation of single readings, not
the standard uncertainty of the five-repeat mean (which would be
$0.3/\sqrt{5} \approx 0.134\,\si{\milli\liter}$). The exercise
asks for the single-reading-spread contribution, so $u_{A} =
0.3\,\si{\milli\liter}$.

\emph{Step 4: combine the contributions in quadrature.} The
resolution and spread are independent (the resolution is a
discretisation effect of the instrument; the spread is from
operator placement and meniscus interpretation). They combine as
\[
u_{c} = \sqrt{u_{\text{res, diff}}^{2} + u_{A}^{2}} =
\sqrt{0.408^{2} + 0.300^{2}} = \sqrt{0.166 + 0.090} \approx
0.506\,\si{\milli\liter}.
\]
Round to one significant figure for the uncertainty: $u_{c}
\approx 0.5\,\si{\milli\liter}$ at $k = 1$.

\emph{Step 5: report.} The volume is $V = 47.5 \pm 0.5\,\si{\milli
\liter}$ at $k = 1$, equivalent to $V = 47.5 \pm 1.0\,\si{\milli
\liter}$ at $k = 2$. Significant figures: the uncertainty falls in
the tenths place, so the reported value is truncated to the same
precision.

\emph{Step 6: the engineering point.} The resolution contribution
($0.408$) dominates the spread contribution ($0.300$) by a factor
of $\sim 1.4$ in standard uncertainty. Reducing the spread (better
technique, more repeats) without improving the resolution would
have diminishing returns: at $u_{A} = 0$ the combined uncertainty
falls only to $0.408\,\si{\milli\liter}$. The right next step is a
finer-resolution graduated cylinder ($0.1\,\si{\milli\liter}$
graduations would drop $u_{\text{res, diff}}$ by $10\times$) or a
mass-based volume measurement using the working density. The
chapter's compatibility-ratio principle applies: if two methods
disagree by more than their combined uncertainty, the
uncertainty model is incomplete.

\repopath{volumes/01-quantity/ch07-length-area-volume-mass/code/displacement\_volume.py}
implements the calculation.

\begin{exercise}
Two independent volume measurements of the same object give
$V_{A} = 47.3 \,\si{\milli\liter}$ with $u_{A} = 1.0 \,\si{\milli
\liter}$ and $V_{B} = 49.1 \,\si{\milli\liter}$ with $u_{B} = 1.2 \,
\si{\milli\liter}$. Compute the compatibility ratio $r$ and state
whether the two methods agree within their combined uncertainty at
the $r \leq 2$ threshold.
\end{exercise}

\paragraph{Solution.}
The compatibility-ratio test for two independent measurements of
the same quantity is the central diagnostic of the chapter's
cross-checking discipline. The solution computes the ratio,
interprets it, and identifies the engineering action.

\emph{Step 1: compute the difference.} The two measurements report
$V_{A} = 47.3$ and $V_{B} = 49.1\,\si{\milli\liter}$, so
\[
\Delta = V_{B} - V_{A} = 1.8\,\si{\milli\liter}.
\]
The sign is positive (method B reports a larger volume), but for
the compatibility test only the magnitude matters.

\emph{Step 2: compute the standard uncertainty of the difference.}
The two methods are independent, so the variances combine in
quadrature:
\[
u(\Delta) = \sqrt{u_{A}^{2} + u_{B}^{2}} = \sqrt{1.0^{2} +
1.2^{2}} = \sqrt{1.00 + 1.44} = \sqrt{2.44} \approx 1.56\,\si{
\milli\liter}.
\]

\emph{Step 3: form the compatibility ratio.} The compatibility
ratio is the normalised difference:
\[
r = \frac{|\Delta|}{u(\Delta)} = \frac{1.8}{1.56} \approx 1.15.
\]

\emph{Step 4: apply the $r \leq 2$ threshold.} Under the working
convention, $r \leq 2$ is taken as compatibility (the two methods
agree within their combined standard uncertainty at the $k = 2$
level). At $r \approx 1.15$, the result is well inside the
compatibility band; no investigation of a previously
uncharacterised error source is triggered.

\emph{Step 5: the sceptical reading.} An $r$ between $1$ and $2$
sits between cleanly-agreeing ($r \ll 1$) and warning ($r > 2$).
For a high-stakes decision (a regulatory submission, a safety-
critical specification), the working laboratory would tighten
either uncertainty before relying on the difference. For the
chapter's volume-measurement example, the next step is a finer-
resolution graduated cylinder for method A and a tighter
calibration on method B. For a low-stakes decision, $r \approx
1.15$ is treated as agreement and the work proceeds.

\emph{Step 6: the engineering point.} The compatibility ratio
encodes the chapter's discipline in one number. A measurement that
agrees with an independent check at $r \ll 2$ is trustworthy; one
that disagrees at $r > 2$ is evidence the uncertainty budget is
incomplete; one that sits in the $1$-to-$2$ band is in the
warning zone and the analyst chooses whether to investigate based
on the decision's stakes. The Hubble case from the chapter's
failure section is the high-stakes example of how an $r > 2$
disagreement was set aside under schedule pressure and how the
chapter's discipline would have demanded otherwise.

\subsection*{Derivation}\pathtag{standard}

\begin{exercise}
Starting from the chapter 4 product-of-powers rule, show that the
relative uncertainty of a sphere's volume $V = \tfrac{4}{3}\pi R^{3}$
satisfies $u(V)/V = 3 u(R)/R$. State the assumption that fails when
the radius is measured at multiple correlated points around the
sphere.
\end{exercise}

\paragraph{Solution.}

\emph{Step 1: start from the volume formula.}
$V = (4/3)\pi R^{3}$.

\emph{Step 2: identify the variables and their exponents.}
The volume is a power-law function of the single variable $R$ with
exponent $+3$. The factor $(4/3)\pi$ is an exact multiplicative
constant; the product-of-powers rule treats it as carrying no
uncertainty.

\emph{Step 3: apply the rule with one variable.}
\[
\left(\frac{u(V)}{V}\right)^{2} = (3)^{2} \left(\frac{u(R)}{R}
\right)^{2}.
\]
Taking the positive square root gives $u(V)/V = 3 \, u(R)/R$, the
required result.

\emph{Step 4: interpret the factor of three.}
A $1\,\%$ relative uncertainty on the radius produces a $3\,\%$
relative uncertainty on the volume. The amplification is the
worked example of the exponent dominating: high-power variables
demand precision investment, because the volume's uncertainty
scales with the radius's relative uncertainty multiplied by the
exponent.

\emph{Step 5: identify the assumption that fails for multipoint
readings.}
The single-variable derivation assumes the volume is a function of
a single number $R$. When the radius is read at $n$ points around
the sphere as $R_{1}, \ldots, R_{n}$ and combined into a mean
$\bar R$, the calculation that propagates only $u(\bar R)$
underestimates the volume's uncertainty whenever the $R_{i}$
share a common systematic error, for example a calibration
offset, a thermal bias, or a parallax convention. The
multi-variable propagation formula carries cross-terms,
\[
\Var(\bar R) = \frac{1}{n^{2}} \sum_{i} \Var(R_{i}) +
\frac{1}{n^{2}} \sum_{i \ne j} \rho_{ij} \, u(R_{i}) \, u(R_{j}),
\]
in which the off-diagonal terms (correlation coefficients
$\rho_{ij}$ across paired readings) survive averaging. For
fully-correlated systematic errors ($\rho_{ij} = 1$ across all
pairs), the variance of $\bar R$ equals the variance of a single
$R_{i}$, and the averaging delivers no precision benefit at all.

\emph{Step 6: the mitigation.}
Calibrate against a reference whose systematic error is
independent of the sphere-measurement chain, by the cross-check
pattern of section 7.5. The cross-check breaks the
common-systematic correlation between readings that the
single-instrument averaging cannot. A useful diagnostic in
practice: if the sample standard deviation of the $R_{i}$ is
suspiciously small ($\ll$ the instrument's stated resolution),
the $R_{i}$ are almost certainly sharing a systematic; the
averaging is reporting precision the chain does not have.

\begin{exercise}
For a triaxial ellipsoid of semi-axes $a, b, c$, the volume is
$V = \tfrac{4}{3}\pi a b c$. Derive the relative standard uncertainty
of $V$ assuming $a, b, c$ are uncorrelated. Identify the case in
which the relative uncertainty of $V$ is dominated by the smallest
semi-axis.
\end{exercise}

\paragraph{Solution sketch.}
Exponents $(+1, +1, +1)$ on the three semi-axes, all uncorrelated.
The product-of-powers rule gives
\[
\left(\frac{u(V)}{V}\right)^{2}
  = \left(\frac{u(a)}{a}\right)^{2}
  + \left(\frac{u(b)}{b}\right)^{2}
  + \left(\frac{u(c)}{c}\right)^{2}.
\]
The smallest semi-axis dominates $u(V)/V$ when its absolute
uncertainty is comparable to the others (a typical case, since the
caliper resolution is shared across all three measurements); then
$u(a)/a, u(b)/b, u(c)/c$ scale inversely with the semi-axis and the
smallest semi-axis contributes the largest relative term. For the
worked example $a = 2.6, b = 2.0, c = 1.8 \,\si{\centi\meter}$ with
$u = 0.1 \,\si{\centi\meter}$ on each, $(u(V)/V)^{2} = 0.0015 +
0.0025 + 0.0031 = 0.0071$, dominated by $c$ at $0.0031$.

\begin{exercise}
Show that, for the difference of two independent volume
measurements $\Delta = V_{A} - V_{B}$ each with standard
uncertainties $u_{A}$ and $u_{B}$, the standard uncertainty of the
difference is $u(\Delta) = \sqrt{u_{A}^{2} + u_{B}^{2}}$. State the
case in which the difference's uncertainty is dominated by the less
precise method, and the case in which both methods contribute
comparably.
\end{exercise}

\paragraph{Solution sketch.}
$\Delta = V_{A} - V_{B}$. By the sum rule for linear combinations
with the coefficient $-1$ on $V_{B}$ and independence, $\Var(\Delta)
= 1^{2} \Var(V_{A}) + (-1)^{2} \Var(V_{B}) = u_{A}^{2} + u_{B}^{2}$.
Hence $u(\Delta) = \sqrt{u_{A}^{2} + u_{B}^{2}}$. The less precise
method dominates whenever $\max(u_{A}, u_{B}) \gg \min(u_{A},
u_{B})$; in that regime $u(\Delta) \approx \max(u_{A}, u_{B})$,
because the smaller variance is negligible in the sum. The methods
contribute comparably when $u_{A} \approx u_{B}$; then $u(\Delta)
\approx \sqrt{2} \, u_{A}$, and tightening either method
proportionally tightens the difference.

\begin{exercise}
Show that, for a balance comparison in which an unknown mass $m$ is
matched against reference masses summing to $m_{\mathrm{ref}}$ with
combined standard uncertainty $u(m_{\mathrm{ref}})$, the standard
uncertainty of $m$ at the balance's null point is bounded below by
$u(m_{\mathrm{ref}})$. Explain why no amount of repeated nulling
reduces the uncertainty of $m$ below the reference's own
uncertainty.
\end{exercise}

\paragraph{Solution sketch.}
At the null point, $m = m_{\mathrm{ref}}$ within the balance's
indication resolution. Treating the balance reading $r$ (the
indicator deflection at the chosen reference loading) as a Type A
random variable with standard uncertainty $u_{\mathrm{bal}}/\sqrt{n}$
after $n$ nulling repeats, the sum rule gives
$u^{2}(m) = u^{2}(m_{\mathrm{ref}}) + u_{\mathrm{bal}}^{2}/n$.
As $n \to \infty$, the balance-noise term vanishes but $u(m) \to
u(m_{\mathrm{ref}})$, which is the asymptotic floor. The
reference's uncertainty is a systematic input to the measurement,
not a random one; averaging cannot reduce a systematic input
(chapter 4 sec.\ on random vs systematic error). To go below the
floor requires either a more accurate reference (a shorter
traceability chain) or a different physical principle (the Kibble
balance vs an artefact-based reference, section 7.3.4).

\subsection*{Estimation}\pathtag{core}

\begin{exercise}
Estimate the volume of water displaced when a typical adult enters
a household bathtub filled to a typical depth. State your
assumptions about the adult's body density and the bathtub's
cross-section. Identify the dominant uncertainty source in the
estimate.
\end{exercise}

\paragraph{Solution sketch.}
Adult mass $\approx 70 \,\si{\kilogram}$, body density close to
water ($\rho \approx 1000 \,\si{\kilogram\per\meter\cubed}$, varying
with body composition by a few percent), so body volume $\approx
70 \,\si{\liter} = 0.070 \,\si{\meter\cubed}$. Bathtub footprint
$\approx 0.7 \times 1.5 \,\si{\meter} = 1.05 \,\si{\meter\squared}$
at the water line. Submerged fraction of the body $\approx 0.5$ to
$0.7$ when reclining at typical water depth, so displaced volume
$\approx 0.04$ to $0.05 \,\si{\meter\cubed}$. Predicted water-line
rise $\approx 0.04 / 1.05 \approx 0.04 \,\si{\meter} = 4 \,\si
{\centi\meter}$. Dominant uncertainty source: the submerged
fraction (a factor-of-1.5 spread depending on posture), then the
bathtub cross-section (a factor-of-1.3 spread across common shapes).
Body density contributes only a few percent. The estimate places
the bathtub-overflow risk for a typical adult at the
several-centimetre level: factor-of-three confidence range is about
$1.5$ to $6 \,\si{\centi\meter}$.

\begin{exercise}
Estimate the surface area of an adult human body, in square metres,
using a body-as-cylinders or body-as-boxes model. State your method
and factor-of-three confidence range. Then compare your estimate to
a cited clinical body-surface-area formula supplied in the solution
notes.
\end{exercise}

\paragraph{Solution sketch.}
Body-as-cylinders model: torso a cylinder of radius $0.15 \,\si
{\meter}$, height $0.6 \,\si{\meter}$, surface area $\approx 2\pi R
H + 2\pi R^{2} \approx 0.71 \,\si{\meter\squared}$. Two arms,
cylinders of radius $0.04 \,\si{\meter}$, length $0.7 \,\si{\meter}$
each: surface area $\approx 2 \times 2\pi R H \approx 0.35 \,\si
{\meter\squared}$. Two legs, cylinders of radius $0.07 \,\si{\meter}$,
length $0.85 \,\si{\meter}$ each: surface area $\approx 2 \times
2\pi R H \approx 0.75 \,\si{\meter\squared}$. Head, sphere of radius
$0.1 \,\si{\meter}$: $4\pi R^{2} \approx 0.13 \,\si{\meter\squared}$.
Total $\approx 1.9 \,\si{\meter\squared}$. The Mosteller formula
$\text{BSA} = \sqrt{H_{\mathrm{cm}} \times M_{\mathrm{kg}} / 3600}$
for a $1.75 \,\si{\meter}$, $70 \,\si{\kilogram}$ adult gives
$\sqrt{175 \times 70 / 3600} \approx 1.84 \,\si{\meter\squared}$.
Factor-of-three confidence range: about $1.0$ to $3.0 \,\si{\meter
\squared}$, comfortably bracketing both estimates.

\begin{exercise}
Estimate the mass of a fully laden road freight vehicle. Use a
stated legal gross combination mass supplied in the problem or
solution notes, then estimate the typical loaded fraction and report
the result with a factor-of-three confidence
range.
\end{exercise}

\paragraph{Solution sketch.}
Stated reference (EU and UK Construction and Use Regulations,
current as of 2025): maximum gross combination mass for a
five-axle articulated lorry is $44 \,\si{\tonne} = 44000 \,\si
{\kilogram}$. Empty tractor and trailer combined $\approx 14 \,\si
{\tonne}$, leaving up to $30 \,\si{\tonne}$ of payload. Typical
loaded fraction varies by cargo: dense cargo (steel, beverages)
fills the weight limit before the volume limit, so the vehicle
runs close to $44 \,\si{\tonne}$; low-density cargo (parcels,
foam, packaged consumer goods) hits the volume limit first, with
gross mass often $20$ to $30 \,\si{\tonne}$. A round estimate for
a representative loaded long-haul lorry is $30 \,\si{\tonne}$ gross,
with factor-of-three confidence range $15$ to $44 \,\si{\tonne}$
(the upper bound capped by the legal limit; the lower bound by
empty-ish backhaul running).

\begin{exercise}
Estimate the density of a single grain of rice and use it to
estimate the volume of one kilogram of rice as packaged for retail
sale. State your assumptions about grain mass and packing density.
Compare your estimate to the volume of a typical $1 \,\si{\kilogram}$
rice package.
\end{exercise}

\paragraph{Solution sketch.}
A long-grain rice kernel is approximately a prolate spheroid of
length $7 \,\si{\milli\meter}$ and diameter $2 \,\si{\milli\meter}$.
Volume $\approx (4/3)\pi (1)(1)(3.5) \approx 15 \,\si{\milli
\meter\cubed}$. A grain mass of about $0.025 \,\si{\gram}$ gives
grain density $\approx 25 \,\si{\milli\gram} / 15 \,\si{\milli\meter
\cubed} \approx 1700 \,\si{\kilogram\per\meter\cubed}$ (consistent
with dried-starch density above water). Random close packing of
elongated grains gives a packing fraction around $0.6$, so the bulk
density of packaged rice is $\approx 1700 \times 0.6 \approx 1000 \,
\si{\kilogram\per\meter\cubed}$. One kilogram of rice therefore
occupies about $1 \,\si{\liter}$ or $1000 \,\si{\milli\liter}$.
Typical $1 \,\si{\kilogram}$ retail rice packages are nominal volume
$1.1$ to $1.3 \,\si{\liter}$ (a settling-and-handling margin), in
agreement with the estimate within $30\,\%$.

\subsection*{Simulation}\pathtag{mastery}

\begin{exercise}
A triaxial ellipsoid has semi-axes $a = 2.6 \,\si{\centi\meter}$,
$b = 2.0 \,\si{\centi\meter}$, and $c = 1.8 \,\si{\centi\meter}$,
each with standard uncertainty $0.1 \,\si{\centi\meter}$. Run or
outline a Monte Carlo simulation with at least $10\,000$ draws to
estimate the distribution of $V = \tfrac{4}{3}\pi a b c$. Compare
the simulated standard deviation with the product-of-powers result
from chapter 4.
\end{exercise}

\paragraph{Solution sketch.}
\repopath{volumes/01-quantity/ch07-length-area-volume-mass/code/ellipsoid\_volume\_mc.py}
implements the simulation. Nominal $V = (4/3)\pi (2.6)(2.0)(1.8)
\approx 39.2 \,\si{\centi\meter\cubed}$. Linear propagation:
$(u(V)/V)^{2} = (0.1/2.6)^{2} + (0.1/2.0)^{2} + (0.1/1.8)^{2}
\approx 0.0071$, giving $u(V)/V \approx 0.084$ and $u(V) \approx
3.3 \,\si{\centi\meter\cubed}$. The Monte Carlo simulation with
$10^{5}$ Gaussian draws on each semi-axis returns a mean $V$
within $0.01\,\%$ of nominal and a sample standard deviation
within a few percent of $3.3 \,\si{\centi\meter\cubed}$. The two
agree because the relative uncertainties are small ($\sim 5\,\%$
each), the linearisation is accurate, and the input distributions
are Gaussian. The Monte Carlo additionally exposes a small
positive skew in $V$ (heavier upper tail) that the linear method
cannot report; for these small relative uncertainties the skew is
ignorable, but at $20\,\%$ relative input uncertainty the
Monte Carlo distribution becomes visibly skewed and the linear
estimate begins to mis-report the typical fluctuation.

\subsection*{Design}\pathtag{standard}

\begin{exercise}
Design a household procedure for measuring the volume of a piece of
fruit too large to fit in any available graduated cylinder. State
the equipment, the procedure, the dominant uncertainty source, and
the expected expanded uncertainty at $k = 2$.
\end{exercise}

\paragraph{Discussion.}
A defensible reply names: (i) an overflow vessel (a wide bowl placed
in a larger catch tray, filled to brimming), (ii) a kitchen scale
with $1 \,\si{\gram}$ resolution to weigh the overflow water rather
than measure its volume directly, (iii) a procedure: tare the catch
tray empty, fill the bowl to overflowing into the catch tray and
discard that water, refill the bowl to overflowing once more,
discard, then place the fruit in the bowl and weigh the overflow in
the catch tray. The volume is the overflow mass divided by water
density at room temperature ($\approx 998 \,\si{\kilogram\per\meter
\cubed}$). Dominant uncertainty source: incomplete brimming of the
bowl (a few millilitres per fill); a secondary source is water
adhering to the fruit on removal. Expected expanded uncertainty at
$k = 2$: about $\pm 5 \,\si{\milli\liter}$ on a $300 \,\si{\milli
\liter}$ fruit volume, or $\pm 2\,\%$. The reply must also state
the failure mode: a fruit with cut surfaces or visible pores will
absorb water and bias the reading low; the procedure works only on
intact, non-porous fruit.

\begin{exercise}
Design a workshop procedure for measuring the volume of an irregular
metal casting whose volume must be known to within $1\,\%$ for a
density-based identification of the alloy. State the instruments,
the procedure, the uncertainty budget, and the dominant uncertainty
source. Identify whether water displacement, geometric approximation,
or image-based reconstruction is the most defensible primary method
for this tolerance.
\end{exercise}

\paragraph{Discussion.}
The $1\,\%$ tolerance on a metal-casting volume effectively
forecloses geometric approximation (typically $5$ to $10\,\%$ error
on a non-trivial casting) and amateur photogrammetry (typically
$2$ to $5\,\%$). Water displacement with a sufficiently fine
overflow technique is defensible: a Mettler-Toledo-class density-
determination kit or an equivalent overflow vessel coupled with a
laboratory balance (resolution $\leq 0.01 \,\si{\gram}$ on a few-
hundred-gram sample) can reach $0.3$ to $0.5\,\%$ volume uncertainty
once temperature, air buoyancy, and water purity are controlled. A
defensible reply lists: the temperature-controlled water bath, the
laboratory balance under controlled environmental conditions, a
sample-degassing step (immersion in a vacuum desiccator to remove
surface bubbles, or sonication in a wetting-agent bath), and a
correction for the water's density at the recorded temperature. The
uncertainty budget includes balance reading ($\sim 0.05\,\%$),
water-density temperature uncertainty ($\sim 0.02\,\%$ per $\si
{\degreeCelsius}$ near room temperature), and a residual surface-
adhesion term ($\sim 0.2\,\%$). The dominant source is the
surface-adhesion term; the failure mode is incomplete wetting,
which biases the measured volume low and the inferred density high.

\begin{exercise}
Design a calibration procedure for a kitchen scale using one or more
low-stakes references: coins of published mass, a sealed package
with stated net mass, or an inexpensive calibration mass set. State
the procedure, the dominant uncertainty source, and the regime in
which the procedure is a verification check rather than a traceable
calibration.
\end{exercise}

\paragraph{Discussion.}
A defensible reply names: (i) reference selection (a sealed
$500 \,\si{\gram}$ bag of dried legumes has net-mass uncertainty
typically $\pm 2 \,\si{\gram}$ from packaging regulations, often
better; an inexpensive M1- or M2-class calibration set has stated
expanded uncertainty under $50 \,\si{\milli\gram}$ on a $1 \,\si
{\kilogram}$ mass at $k = 2$), (ii) the procedure: weigh the
reference five times, removing and re-centring between weighings,
and report the mean and sample standard deviation; tabulate the
indication-versus-reference relation across at least three points
spanning the scale's range, (iii) the dominant uncertainty source:
the reference's stated uncertainty if a calibration set, the
package-net-mass tolerance if a retail item. The procedure is a
verification check, not a traceable calibration, because the
references are not themselves linked through an unbroken
documented chain to an SI realisation; for a traceable calibration
the reference must come with a calibration certificate from an
ISO/IEC 17025-accredited laboratory \cite{std:iso-iec-17025}.

\begin{exercise}
Design a measurement plan for the surface roughness of a workshop
benchtop, given access only to household instruments (a smartphone
camera, a sheet of fine paper, a flat reference like a kitchen
worktop edge, a flashlight). State the qualitative comparison the
plan can deliver and identify what a profilometer would add that the
household plan cannot.
\end{exercise}

\paragraph{Discussion.}
A defensible household plan delivers a qualitative ordering, not a
quantitative $R_{a}$. The reply names: (i) a grazing-light test
(flashlight held nearly parallel to the surface, photographed by
the smartphone; surface ridges cast long shadows whose contrast and
length correlate with roughness), (ii) a finger-touch comparison
against a known-smoother reference (the kitchen worktop edge or a
flat ceramic plate), (iii) a paper-pull test (a sheet of fine paper
slid across the surface produces a faint scratch or none, depending
on grit). The plan distinguishes "as-machined" from "polished" but
cannot give $R_{a}$ in micrometres. A profilometer adds: a
quantitative $R_{a}$ value tied to a standard, repeatable
measurement over a defined evaluation length and a defined cutoff
or nesting index per ISO 21920-2 \cite{std:iso-21920-2-2021}, and
the ability to compare two surfaces whose qualitative grazing-light
appearance is indistinguishable. The household plan is a
screening tool; the profilometer is a measurement instrument.

\subsection*{Diagnosis and reverse engineering}\pathtag{core}

\begin{exercise}
A water-displacement measurement of a stone gives a volume that
varies by $\pm 5\,\%$ between repeated readings. State three plausible
mechanisms and propose the next measurement that distinguishes them.
\end{exercise}

\paragraph{Solution.}

\emph{Step 1: state the observed signature.}
The water-displacement reading varies by $\pm 5\,\%$ across
repeats on a single stone, an unusually wide spread for a method
whose graduated-cylinder resolution-derived uncertainty is
typically a fraction of a percent on a $\sim 100 \,\si{\milli
\liter}$ object. The width is a signal that one or more
unmodelled error sources are active.

\emph{Step 2: enumerate plausible mechanisms.}
(i) Trapped air bubbles released differently between readings.
Each released bubble shifts the displaced volume downward in a
discrete step, producing a non-Gaussian distribution of readings
clustered around two or three levels rather than a smooth spread.
(ii) Parallax in reading the meniscus: a Type A spread that should
average toward zero over many readings, with a characteristic
magnitude set by the graduated cylinder's resolution and the
operator's eye position.
(iii) Absorption by the stone: a porous sandstone or a stone with a
crack network soaks up a small mass of water on each immersion,
biasing the subsequent reading downward by the absorbed volume.
The bias accumulates if the stone is not re-dried between
readings.

\emph{Step 3: design a test that distinguishes the mechanisms.}
Pre-soak the stone in water for an hour, then dry the surface (a
quick blot, not a thorough dry that would re-expose dry pores),
then repeat the five readings under tap-gentle conditions to
prevent new bubble capture. The pre-soak saturates any absorbent
pore network; subsequent readings should be free of the
absorption-induced shift.

\emph{Step 4: read off the diagnosis from the new spread.}
If the spread collapses to a fraction of a percent, bubbles or
absorption were dominant. If the spread persists at $\pm 5\,\%$,
parallax (or another reading-spread source) is dominant. The
distinction between bubbles and absorption requires one more
measurement: weigh the stone before and after the pre-soak. An
absorbent stone gains measurable mass on the order of $0.1$ to
$1 \,\si{\gram}$, depending on porosity; a non-absorbent stone
gains nothing within the kitchen scale's resolution.

\emph{Step 5: prescribe the fix for the dominant mechanism.}
For bubbles: degas the water by light boiling and slow cooling,
or by ultrasonic agitation; immerse the stone slowly and tap the
cylinder wall to dislodge surface bubbles. For absorption: report
two volumes (saturated and dry), and use the saturated volume for
the displacement-derived density calculation. For parallax: replace
the graduated cylinder with a finer-graduated cylinder, or take a
digital photograph of the meniscus from a fixed perpendicular
viewpoint, with the cylinder backlit through a millimetre grid.

\emph{Step 6: capture the discipline.}
A wide-spread water-displacement reading is the chapter 4
fingerprint of an unmodelled error source. The diagnosis runs in
sequence: characterise the distribution of the readings; propose
mechanisms that fit the distribution's shape; design a test that
distinguishes them; localise; fix. The pattern is independent of
the measurand; the volume-of-a-stone case is a tutorial example
of the chapter 7 cross-checking discipline applied within a
single measurement chain.

\begin{exercise}
A density measurement of a metal sample gives
$8.95 \,\si{\gram\per\centi\meter\cubed}$, suggesting the sample is
copper, but the sample is suspected to be a copper-zinc brass.
State which input (mass or volume) is most likely the source of the
discrepancy and what additional measurement would resolve the
ambiguity.
\end{exercise}

\paragraph{Solution sketch.}
Brass densities range from about $8.4$ to $8.7 \,\si{\gram\per
\centi\meter\cubed}$ for typical Cu-Zn compositions; pure copper is
$8.96 \,\si{\gram\per\centi\meter\cubed}$. The discrepancy is at
most $5\,\%$ and most likely from the volume input: a mass
measurement on a small metal sample with an analytical balance is
typically $\sim 0.01\,\%$, while a water-displacement or
geometric-approximation volume is typically $1$ to $5\,\%$. The
candidate volume mechanisms are: trapped air bubbles biasing
volume low and density high, or geometric-approximation model
mismatch. The resolving measurement is an independent volume
determination by a higher-precision method: a hydrostatic balance
(measuring the sample's apparent weight loss in water, a direct
volume-by-buoyancy method with $0.1$ to $0.3\,\%$ uncertainty) is
the natural choice. An independent identification path is an X-ray
fluorescence (XRF) measurement of the sample's elemental
composition; XRF resolves copper from brass directly without
relying on density.

\begin{exercise}
A bathroom scale calibrated at the manufacturer's facility in
northern Europe consistently reads about $0.5\,\%$ low after the
scale is shipped to a customer in Quito. Identify the mechanism, and
state whether the issue is a calibration error, a mechanical
failure, or an inherent limitation of the instrument's principle.
\end{exercise}

\paragraph{Solution.}

\emph{Step 1: name the measurement principle of the instrument.}
A force-measuring bathroom scale uses a strain-gauge load cell. A
mass $m$ on the pan produces a weight $W = m g$ on the cell; the
cell's strain gauges report a voltage proportional to $W$; the
electronics divide by an assumed value of $g$ (the calibration
site's value) and display a mass. The instrument is not weighing
mass directly; it is reading weight and converting to a
mass-equivalent through an assumption about gravity.

\emph{Step 2: compute $g$ at each location.}
The standard model for $g$ on Earth's surface combines latitude
and altitude. At sea level the gravitational acceleration varies
from $\approx 9.780 \,\si{\meter\per\second\squared}$ at the
equator to $\approx 9.832 \,\si{\meter\per\second\squared}$ at the
poles. The international gravity formula (the 1980 GRS80 form)
gives $g_{0}(\phi) \approx 9.780327 [1 + 0.0053024 \sin^{2}\phi -
0.0000058 \sin^{2}(2\phi)] \,\si{\meter\per\second\squared}$, with
$\phi$ the geographic latitude. At $52^{\circ}$ N (northern
Europe, e.g.\ Berlin, the Netherlands), $g_{0} \approx 9.812 \,\si
{\meter\per\second\squared}$ at sea level. The free-air correction
for altitude is $-3.086 \times 10^{-6} \,\si{\meter\per\second
\squared\per\meter}$.

\emph{Step 3: apply the model to Quito.}
Quito sits at approximately $0.2^{\circ}$ S latitude (call it
$0^{\circ}$ for this computation) and $2850 \,\si{\meter}$ above
sea level. The latitude part: $g_{0}(0) = 9.780 \,\si{\meter\per
\second\squared}$. The altitude part: $-3.086 \times 10^{-6}
\times 2850 = -8.8 \times 10^{-3} \,\si{\meter\per\second\squared}$.
Net: $g \approx 9.780 - 0.009 = 9.771 \,\si{\meter\per\second
\squared}$, a value consistent with published Quito gravity
surveys at $9.769$ to $9.771 \,\si{\meter\per\second\squared}$
depending on the local geological context.

\emph{Step 4: compute the indicated-mass shift.}
The ratio of indicated mass at Quito to indicated mass at the
calibration site, for the same physical mass on the pan, equals
the ratio of $g$ values:
\[
\frac{m_{\mathrm{Quito}}^{\mathrm{indicated}}}
{m_{\mathrm{NorthEU}}^{\mathrm{indicated}}}
  = \frac{g_{\mathrm{Quito}}}{g_{\mathrm{NorthEU}}}
  = \frac{9.769}{9.812} \approx 0.9956.
\]
The scale reads $0.44\,\%$ low in Quito, matching the observed
$\approx 0.5\,\%$ figure within the measurement precision of a
domestic scale.

\emph{Step 5: classify the issue.}
The phenomenon is an inherent limitation of the force-measuring
principle. The instrument was calibrated correctly for the
calibration site, and is operating to its specifications at the
new site. It is not a calibration error in the manufacturing
sense; it is not a mechanical failure of the load cell. The
phenomenon is the expected behaviour of a force-measuring
instrument moved to a location with different $g$. The
manufacturer's datasheet, if comprehensive, will state the
calibration site and the expected geographic correction; in
practice most domestic-scale datasheets do not, and the customer
attributes the shift to a defect.

\emph{Step 6: the fix and the principle to remember.}
The fix is a local recalibration against a known mass at the new
site: place a calibration mass of known value $m_{\mathrm{ref}}$
on the pan, and adjust the scale's internal gain factor (or apply
a multiplicative correction in the user's records) so the
indication reads $m_{\mathrm{ref}}$. A balance comparing mass
against mass (the equal-arm pattern of
\cref{fig:vol01:ch07:equal-arm-balance}, or a force-restoration
balance calibrated against an internal mass standard) would not
have shown the shift, because $g$ cancels in the comparison. The
deeper lesson is geometric: the structural cancellation in the
equal-arm pattern is the reason mass-against-mass comparison was
the metrological standard for centuries before electronic load
cells, and it remains the standard wherever the measurement must
travel.

\begin{exercise}
A geometric approximation of a stone's volume by an ellipsoid model
gives a result $20\,\%$ smaller than the water-displacement
measurement. State two plausible mechanisms (the model is too
smooth, the displacement reading included surface bubbles) and
propose, for each, the test that would confirm or refute it.
\end{exercise}

\paragraph{Solution sketch.}
The geometric ellipsoid is smaller than the displacement reading
when (a) the actual stone has ridges and bulges the ellipsoid
treats as smooth (model underestimate) and (b) the displacement
reading is biased high because trapped air bubbles displace water
additional to the stone (reading overestimate). Both push the ratio
$V_{\mathrm{displacement}} / V_{\mathrm{geometric}}$ above $1$.
Test for (a): take additional caliper readings at the bulges and
re-fit a more refined model (a sum of simple shapes rather than a
single ellipsoid); if the geometric estimate rises toward the
displacement, the original model was the dominant source. Test for
(b): pre-soak the stone for an hour, then re-do the displacement
under tap-gentle conditions. If the displacement falls toward the
geometric estimate, trapped bubbles were dominant. The most likely
finding for a typical garden stone is that both mechanisms
contribute; the diagnostic is which dominates, and the answer
determines whether the project's $20\,\%$ disagreement is closed
by remodelling the geometry or by remeasuring the displacement.

\subsection*{Failure analysis}\pathtag{core}

\begin{exercise}
The Hubble Space Telescope's primary-mirror figure was measured
against a reflective null corrector that was itself wrong. Argue, in
200 to 300 words, why the auxiliary tests (the refractive null
corrector and the inverse null test) were the wrong-ruler detector
that the process needed, and identify the procedural barrier whose
absence allowed the auxiliary tests' indications to be set aside.
Cite the OSBI report \cite{acc:nasa-hubble-optical-systems-1990}.
\end{exercise}

\paragraph{Solution.}

\emph{Step 1: state the chain's internal consistency.}
The reflective null corrector was internally consistent. Its
assembly chain ran: a metering rod whose features defined the
axial spacings, against which a field lens was located; the
located lens produced the RNC's reference figure; the polished
mirror was figured against that reference. Every link reported
agreement with the link above it. The chain was internally
auditable.

\emph{Step 2: identify what internal consistency could not catch.}
What no internal check could detect was that the metering-rod
reading was misinterpreted at the top of the chain. The rod had
two features near its end (an end cap and a target feature,
recessed by $\sim 1.3 \,\si{\milli\meter}$) that the assembler
could mistake for one another. The lens was set against the end
cap rather than the target feature, the RNC was thereafter
internally consistent with the wrong axial spacing, and the mirror
was polished to match
\cite{acc:nasa-hubble-optical-systems-1990}. An internally
consistent chain is one whose links agree with each other; it is
not one whose top is set against the right reference.

\emph{Step 3: name the auxiliary tests as a different chain.}
The auxiliary tests, a refractive null corrector and an inverse
null test, were a different measurement chain. They used different
physical principles (refractive optics rather than reflective),
different instruments (a separate refractive corrector, an
inverse-null configuration that did not rely on the RNC's
geometric definition), and different sources of systematic error.
Their independence from the RNC chain was precisely what made them
useful as cross-checks.

\emph{Step 4: quantify the disagreement they reported.}
The auxiliary tests' indications at fabrication time disagreed
with the RNC's reading by an amount inconsistent with their
combined uncertainty budgets. The disagreement was not noise; it
was the compatibility-ratio signal of section 7.5 arriving in
real time. A working metrologist running the chapter 4
compatibility ratio on the two reports would have flagged the
inconsistency, and would have classified it as requiring
resolution before the mirror entered service.

\emph{Step 5: name the procedural barrier whose absence
mattered.}
The barrier was a formal anomaly-resolution requirement that the
auxiliary-test indication be treated as a finding requiring
resolution before proceeding, rather than as advisory information
that could be overruled by the more authoritative instrument under
cost and schedule pressure. The Allen Commission's Chapter 9
recommendations \cite{acc:nasa-hubble-optical-systems-1990} are
that requirement, framed in optical-fabrication terms: anomaly
resolution must be procedural, must include the customer (NASA)
in the decision chain, and must not allow informal seniority of
one instrument over another to determine which test is believed.

\emph{Step 6: state the lesson for length measurement at every
scale.}
The Hubble case is the wrong-ruler pattern at the largest optical
scale the engineering record contains. A reference instrument is
itself a measurement chain; the cross-check on the reference is an
independent measurement of the reference, not a repeat of the same
measurement on the same chain. The discipline scales down without
modification to the workshop: a calliper whose calibration is
verified against a gauge block from the same calibration lot,
returned from the same laboratory, is internally consistent and
externally unverified.

\begin{exercise}
A length measurement chain is documented in three written links: a
calliper calibrated against a gauge block, a gauge block calibrated
against a master reference, the master reference calibrated against
an SI realisation. State the failure mode that occurs when the
master reference is itself drifting between calibration cycles, and
the cross-check that would detect the drift before the next external
recalibration.
\end{exercise}

\paragraph{Solution.}

\emph{Step 1: walk the chain forward to find the failure mode.}
The calliper is calibrated against the gauge block; the gauge
block is calibrated against the master reference; the master
reference is calibrated against an SI realisation, but only at
the formal recalibration interval. Between recalibrations the
master drifts (handling wear, contamination, environmental
exposure). Each downstream calibration step calibrates against
the master at its present, drifted state, so the entire chain
reports lengths shifted by the master's drift. The reports are
internally self-consistent, because every link agrees with the
link above it; they are externally biased, because the top of the
chain has moved.

\emph{Step 2: identify why no internal check catches it.}
The chain's internal-consistency check is each link's agreement
with the link above. The check is structurally blind to a drift
at the top, because it has no second top to compare against. The
calliper's calibration certificate, the gauge block's calibration
certificate, and the master's calibration certificate all match
their stated reference values, because the references are the
same drifting object measured from the same chain.

\emph{Step 3: name the cross-check that detects the drift.}
An independent reference of comparable accuracy, held outside the
same calibration laboratory: a second master reference from a
different supplier or a different lot, calibrated through a
different traceability chain. The two masters are then periodically
compared within the laboratory (a within-laboratory
cross-comparison) by mounting both on the same comparator and
running a calibration-grade differential measurement.

\emph{Step 4: specify the test.}
Run the compatibility-ratio test of section 7.5 on the two
masters: $\Delta = L_{1} - L_{2}$, $u(\Delta) = \sqrt{u(L_{1})^{2}
+ u(L_{2})^{2}}$, $r = |\Delta| / u(\Delta)$. As long as $r < 2$,
the two masters agree within their combined uncertainty; the
chain reports nothing actionable. When $r$ crosses $2$, the
working assumption that the chain's reference is stable has
broken; the appropriate response is to schedule an external
recalibration of both masters before further work issues from
the chain.

\emph{Step 5: localise which master is drifting.}
The compatibility-ratio test tells the reader the two masters
disagree; it does not tell the reader which has drifted. The
localisation step is a third independent reference, or a return
of both to the external calibration laboratory, or a fixed
external comparison against an SI realisation (a different
length-measurement primary chain, an interferometer tied to a
stabilised frequency). The general pattern: $n = 2$ references
catch a drift; $n = 3$ references localise it; $n \geq 4$
references begin to characterise the drift's statistical
properties.

\emph{Step 6: generalise to the IPK lesson.}
The discipline restates the IPK lesson of section 7.3.5 at the
working laboratory scale. A single artefact cannot detect its own
drift; a cohort of nominally identical artefacts can, by
within-cohort comparison. The IPK case operated this principle in
the negative: the cohort showed scatter, but the chain's
definitional structure prevented declaring any one artefact
drifted relative to the others. The working laboratory recovers
the affirmative version by holding multiple masters through
distinct calibration chains, and by treating the cross-check as
the routine check that the chain's top is intact.

\begin{exercise}
A workshop reports that a part dimensioned in millimetres on its
drawing has been manufactured against a CMM whose internal frame
was rotated $0.5^{\circ}$ relative to the workpiece's design
frame. State the resulting bias on a $300 \,\si{\milli\meter}$
length measured along the design x-axis, and the audit step at which
the rotation should have been caught.
\end{exercise}

\paragraph{Solution sketch.}
A length $L$ measured along the design x-axis but reported in a
frame rotated by $\theta = 0.5^{\circ}$ has projected
length $L \cos\theta = L (1 - \theta^{2}/2 + \ldots)$. For
$\theta = 0.5^{\circ} \approx 0.00873$ radian, $\cos\theta \approx
1 - 3.8 \times 10^{-5}$. The reported $x$-component is therefore
$300 \times (1 - 3.8 \times 10^{-5}) \approx 299.989 \,\si
{\milli\meter}$, a bias of about $11 \,\si{\micro\meter}$ low on
a $300 \,\si{\milli\meter}$ length. (The bias is small because
$\cos\theta$ is near unity for small angles, but the orthogonal
$y$-projection of $L \sin\theta \approx 2.6 \,\si{\milli\meter}$
is the larger effect on the part's shape; the workpiece's
diagonal is rotated by half a degree relative to design.) The
audit step that should have caught the rotation is the
\engterm{workpiece alignment} step that defines the
machine-frame-to-design-frame transformation at the start of the
CMM session. The standard practice is a three-point alignment
procedure that explicitly establishes the design frame's axes
against features on the workpiece; the alignment's residual error
should be recorded and propagated into the measurement uncertainty
budget. A workshop that omits or under-reports the alignment step
ships parts whose dimensional history is consistent with the
machine frame rather than the design frame.

\subsection*{Judgment}\pathtag{standard}

\begin{exercise}
Argue, in one paragraph, whether a single high-precision measurement
or two independent moderate-precision measurements give the engineer
more confidence in the value of a quantity. Take a position; defend
it with a specific scenario in which your preferred answer fails.
\end{exercise}

\paragraph{Discussion.}
The defensible position is that two independent moderate-precision
measurements give more confidence in the value of a quantity than a
single high-precision measurement, because the two reports allow a
compatibility-ratio check that the single report cannot. The
single high-precision measurement reports its self-derived
uncertainty, which is bounded by the error sources the operator
knows about; an uncharacterised systematic error inflates the
reported value with no in-band indicator. The two moderate
measurements, by contrast, surface any error source one method
sees and the other does not. The position fails when the two
moderate methods share a dominant common-mode systematic (a shared
temperature reference, a shared calibration laboratory, a shared
operator bias); then the two reports agree because they are wrong
together, and the single high-precision report is more honest.
A defensible reply names the common-mode failure mode explicitly,
because the resolution is then design of genuinely independent
methods rather than addition of a second method drawn from the
same chain.

\begin{exercise}
A senior colleague tells the reader that ``cross-checking is for
graduate students; in industry we trust the calibrated instrument
and move on.'' Write a one-paragraph reply that takes the colleague
seriously, identifies the regime in which the claim is defensible,
and identifies the regime in which it is not.
\end{exercise}

\paragraph{Discussion.}
A defensible reply concedes the colleague's point in its proper
regime: for routine production measurements within a calibrated
instrument's stated working range, on a measurand whose
characteristics match the instrument's calibration model, the
formal cross-check has been done once at certification time and
the per-measurement repeat is wasted motion. The discipline of
calibration is precisely what lets the production engineer skip
the per-measurement cross-check. The reply then identifies the
regime in which the claim is not defensible: any measurement that
matters more than the calibration certificate's expanded
uncertainty allows (a borderline pass/fail dimension on a critical
part), any measurement outside the instrument's calibrated range
(extrapolation), any measurement whose physical principle puts it
in a regime the calibration model did not test (a thermal regime
outside calibration limits, a vibration regime the calibration
laboratory did not simulate), and any measurement made on an
instrument whose calibration certificate has expired. In those
regimes, the cross-check is the discipline; outside them, the
calibration is the cross-check.

\begin{exercise}
Argue, in 200 to 300 words, whether the four geometric quantities of
this chapter (length, area, volume, mass) deserve their privileged
status in the engineering record, or whether other quantities
(temperature, pressure, time) deserve equal status. Take a position
and defend it with reference to the working measurements of a chosen
engineering domain.
\end{exercise}

\paragraph{Discussion.}
A defensible reply takes a position rather than splitting the
difference. One defensible position: the four chapter quantities
are not privileged by their physics but by the engineering
workflow that begins with the artefact's geometry and mass
before its operating conditions are imposed. A drawing is
dimensioned; a build of materials is massed; a quotation is
priced by mass and volume. Temperature, pressure, and time enter
the engineering record only when the artefact is in use, while
length and mass enter at the design stage and remain throughout
its life. The defence: a chosen domain (say, bridge construction)
specifies length tolerances, volume of concrete poured, mass of
steel rebar, and area of formwork before any thermal or temporal
condition is documented; the as-built record reaches the
maintenance phase with the four geometric quantities as the
primary identifying record. The opposing position is also
defensible: in dynamic systems (control, aerospace, power
electronics), time and frequency are the most-measured quantities,
and a high-frequency systems engineer might list them as the
discipline's first quantities. The judgment-grade answer is a
defensible position with a specific domain and a specific
counterexample; the indefensible answer is the both-and reply
that names no working measurement.
